\documentclass[11pt,a4paper]{article}
\usepackage[utf8]{inputenc}
\usepackage[T1]{fontenc}
\usepackage[spanish,es-noshorthands]{babel}
\usepackage[margin=2.5cm]{geometry}
\usepackage{amsmath,amssymb,amsthm,mathtools}
\usepackage{tikz}
\usepackage{pgfplots}
\pgfplotsset{compat=1.18}
\usepackage{booktabs}
\usepackage{array}
\usepackage{enumitem}
\usepackage{hyperref}
\hypersetup{colorlinks=true, linkcolor=blue, citecolor=blue, urlcolor=blue}

\theoremstyle{plain}
\newtheorem{teorema}{Teorema}[section]
\newtheorem{propiedad}[teorema]{Propiedad}
\newtheorem{corolario}[teorema]{Corolario}
\newtheorem{lema}[teorema]{Lema}
\theoremstyle{definition}
\newtheorem{definicion}[teorema]{Definición}
\newtheorem{ejemplo}[teorema]{Ejemplo}
\theoremstyle{remark}
\newtheorem{observacion}[teorema]{Observación}

\renewcommand{\proofname}{Demostración}

\title{Unidad 7: Prueba de hipótesis \\ \large Apunte de teoría}
\author{Probabilidad y Estadística (5765) \\ Universidad Nacional del Sur}
\date{}

\begin{document}
\maketitle
\tableofcontents
\newpage

\section{Conceptos fundamentales de la prueba de hipótesis}

\subsection{De la estimación a la decisión estadística}

En la Unidad 6 desarrollamos la teoría de la \textbf{estimación} de parámetros poblacionales: dado un parámetro desconocido $\theta$ (una media $\mu$, una varianza $\sigma^2$, una proporción $p$, etc.), construimos estimadores puntuales y, más informativamente, intervalos de confianza que acompañan la estimación con una medida de precisión. El problema que abordamos ahora es de naturaleza distinta. En muchas aplicaciones el investigador no necesita estimar $\theta$: ya cuenta con un valor de referencia, una conjetura, una especificación técnica o una afirmación que debe verificarse, y lo que necesita es \textbf{decidir}, a partir de la evidencia que proporciona una muestra aleatoria, si esa conjetura es sostenible o si debe rechazarse en favor de una alternativa.

Este tipo de problema aparece constantemente en la práctica:
\begin{itemize}[nosep]
\item ¿El contenido medio de las bolsas de un producto es realmente el que indica el rótulo?
\item ¿Un nuevo proceso de fabricación reduce la proporción de piezas defectuosas respecto del proceso anterior?
\item ¿La variabilidad de un instrumento de medición supera la tolerancia técnica admitida?
\item ¿El rendimiento académico de los estudiantes es independiente del turno en el que cursan?
\item ¿Dos formulaciones de un medicamento producen, en promedio, el mismo efecto terapéutico?
\end{itemize}

En todos los casos hay una pregunta binaria de fondo (la afirmación es cierta o no lo es) y una fuente de información parcial y aleatoria (la muestra) con la cual formular una respuesta. La disciplina que organiza este tipo de razonamiento se llama \textbf{prueba de hipótesis} (o contraste de hipótesis), y constituye, junto con la estimación, uno de los dos grandes pilares de la inferencia estadística clásica. De hecho, como veremos en la Sección \ref{sec:dualidad}, ambos pilares están profundamente relacionados: toda prueba de hipótesis bilateral puede reformularse en términos de un intervalo de confianza, y viceversa.

Conviene remarcar desde el inicio la naturaleza \emph{probabilística} —no determinística— de las conclusiones que se obtienen. Una prueba de hipótesis nunca demuestra con certeza absoluta que una afirmación sobre la población es verdadera o falsa: lo que hace es cuantificar, mediante probabilidades controladas de antemano, el riesgo de equivocarse al tomar una decisión basada en una muestra finita. Este control explícito del error es lo que distingue a la inferencia estadística de la mera opinión basada en datos.

\subsection{Hipótesis nula y alternativa}

\begin{definicion}
Una \textbf{hipótesis estadística} es una afirmación (conjetura) acerca del valor de uno o más parámetros de una o varias poblaciones, o acerca de la forma de la distribución de una variable aleatoria.
\end{definicion}

Toda prueba de hipótesis involucra dos afirmaciones complementarias.

\begin{definicion}
La \textbf{hipótesis nula}, denotada $H_0$, es la afirmación que se somete a prueba. Representa la situación de referencia, el \emph{status quo}, el valor histórico, la especificación técnica, o la ausencia de efecto o diferencia. Se formula siempre como una igualdad respecto de un valor de referencia $\theta_0$:
\[
H_0:\ \theta=\theta_0.
\]
\end{definicion}

\begin{definicion}
La \textbf{hipótesis alternativa}, denotada $H_1$ (o $H_a$), es la afirmación que se sostiene como razonable cuando la evidencia muestral resulta incompatible con $H_0$. Según el problema, puede formularse de tres maneras:
\[
H_1:\ \theta\neq\theta_0 \quad \text{(\emph{bilateral} o de dos colas)},
\]
\[
H_1:\ \theta>\theta_0 \qquad \text{o bien} \qquad H_1:\ \theta<\theta_0 \quad \text{(\emph{unilaterales} o de una cola)}.
\]
\end{definicion}

La elección entre una alternativa bilateral y una unilateral no es arbitraria: depende exclusivamente de la pregunta de investigación planteada \emph{antes} de observar los datos. Si el interés es detectar cualquier apartamiento del valor de referencia, sin importar el sentido, corresponde una alternativa bilateral. Si el interés específico es detectar un aumento (o una disminución), corresponde una alternativa unilateral en el sentido correspondiente. Decidir el tipo de alternativa después de observar los datos —por ejemplo, mirar el signo de $\overline{x}-\mu_0$ y recién ahí elegir la cola— invalida por completo el procedimiento, porque distorsiona la probabilidad de error que se pretende controlar.

\begin{observacion}
$H_0$ y $H_1$ deben ser mutuamente excluyentes y, en el marco que se utiliza en este curso, exhaustivas respecto del parámetro: toda decisión posible (rechazar o no rechazar $H_0$) se interpreta en términos de estas dos afirmaciones. Es habitual que la afirmación que el investigador sospecha o desea sustentar con evidencia se formule como $H_1$, dejando en $H_0$ la postura conservadora o la que representa ausencia de novedad. Esta convención no es casual: como se verá en la Sección \ref{sec:errores}, la estructura de la prueba está diseñada para controlar rigurosamente el error de rechazar $H_0$ de manera indebida, por lo que conviene que $H_0$ sea la afirmación cuyo rechazo erróneo tenga las consecuencias más costosas.
\end{observacion}

\begin{ejemplo}
Una fábrica de rodamientos afirma que el diámetro medio de su producción es $\mu_0=10$ mm. Un control de calidad sospecha que el proceso se desvió de esa especificación, sin tener razones para pensar que el desvío sea en un sentido particular. Se plantea entonces una prueba bilateral:
\[
H_0:\mu=10 \quad \text{vs.}\quad H_1:\mu\neq10.
\]
\end{ejemplo}

\begin{ejemplo}
Un nuevo fertilizante se ensaya con la esperanza de que \textbf{aumente} el rendimiento medio $\mu$ de un cultivo respecto del valor histórico $\mu_0=40$ qq/ha obtenido con el fertilizante tradicional. La pregunta de investigación tiene un sentido definido (aumento), por lo que corresponde una prueba unilateral:
\[
H_0:\mu\le40 \quad \text{vs.}\quad H_1:\mu>40.
\]
En la práctica, para el cálculo del estadístico y de la región crítica se trabaja siempre con la igualdad $H_0:\mu=40$, que es el punto de la frontera de $H_0$ más desfavorable para el investigador (el más difícil de rechazar), y por lo tanto el que garantiza que la probabilidad de error tipo I no supere el nivel fijado en ningún punto de $H_0:\mu\le 40$.
\end{ejemplo}

\subsection{Errores de tipo I y tipo II}
\label{sec:errores}

Una prueba de hipótesis conduce, inevitablemente, a una de dos decisiones: \textbf{rechazar} $H_0$ o \textbf{no rechazar} $H_0$. Puesto que la decisión se basa en una muestra aleatoria y no en la población completa, siempre existe una probabilidad positiva de que la decisión tomada no coincida con el verdadero estado de la naturaleza. La siguiente tabla resume las cuatro situaciones posibles.

\begin{table}[h]
\centering
\begin{tabular}{@{}lcc@{}}
\toprule
 & $H_0$ verdadera & $H_0$ falsa \\
\midrule
\textbf{No rechazar $H_0$} & Decisión correcta ($1-\alpha$) & Error de tipo II ($\beta$) \\
\textbf{Rechazar $H_0$} & Error de tipo I ($\alpha$) & Decisión correcta ($1-\beta$, potencia) \\
\bottomrule
\end{tabular}
\caption{Errores posibles en una prueba de hipótesis.}
\end{table}

\begin{definicion}
El \textbf{error de tipo I} consiste en rechazar $H_0$ cuando en realidad es verdadera. Su probabilidad se llama \textbf{nivel de significación} de la prueba y se denota
\[
\alpha = P(\text{Rechazar } H_0 \mid H_0 \text{ verdadera}).
\]
\end{definicion}

\begin{definicion}
El \textbf{error de tipo II} consiste en no rechazar $H_0$ cuando en realidad es falsa. Su probabilidad se denota
\[
\beta = P(\text{No rechazar } H_0 \mid H_0 \text{ falsa}).
\]
\end{definicion}

El nivel de significación $\alpha$ es fijado \emph{de antemano} por quien diseña la prueba, típicamente en valores como $0.10$, $0.05$ o $0.01$, y representa la probabilidad máxima que se está dispuesto a tolerar de rechazar erróneamente una hipótesis nula verdadera. Es importante remarcar que $\alpha$ no es una cantidad que se calcule a partir de los datos: es una decisión de diseño, tomada antes de observar la muestra, análoga a fijar el nivel de confianza $1-\alpha$ de un intervalo de confianza.

A diferencia de $\alpha$, la probabilidad $\beta$ de error de tipo II no es un único número: depende de \emph{cuál} sea el verdadero valor del parámetro dentro de $H_1$ (es decir, de qué tan alejado esté el valor real $\theta$ de $\theta_0$). Cuanto más cercano esté el valor verdadero al valor hipotetizado $\theta_0$, más difícil es distinguirlo con una muestra finita, y por lo tanto mayor será $\beta$ para ese valor particular.

\begin{propiedad}[Relación entre $\alpha$ y $\beta$]
Para un tamaño de muestra $n$ fijo y un estadístico de prueba dado, $\alpha$ y $\beta$ están inversamente relacionados: si se reduce el tamaño de la región crítica para disminuir $\alpha$, aumenta correspondientemente $\beta$ (y viceversa). La única manera de disminuir \emph{simultáneamente} ambas probabilidades de error es aumentar el tamaño de la muestra $n$.
\end{propiedad}

Esta propiedad tiene una justificación intuitiva sencilla: si $RR$ denota la región crítica y $RA=RR^c$ la región de no rechazo, entonces $\alpha=P(T\in RR\mid H_0)$ y $\beta=P(T\in RA\mid H_1)$. Si se agranda $RR$ para hacer $\alpha$ más chico, automáticamente $RA$ se agranda también (es su complemento), y por lo tanto $\beta=P(T\in RA\mid H_1)$ tiende a crecer. La única forma de reducir el tamaño de ambas regiones de error sin que una compense a la otra es hacer que la distribución del estadístico bajo $H_0$ y bajo $H_1$ se \emph{concentren} más y se separen más entre sí, lo cual se logra reduciendo la varianza del estimador, es decir, aumentando $n$.

\begin{observacion}[¿Por qué se prioriza el control de $\alpha$?]
En la construcción clásica de una prueba de hipótesis, $H_0$ suele representar la afirmación conservadora, aquella cuyo rechazo erróneo acarrea las consecuencias más costosas (por ejemplo, retirar del mercado un lote de producto en buen estado, interrumpir un proceso de fabricación que funciona correctamente, o acusar de un efecto que no existe). Por esta razón, el procedimiento controla \emph{directamente} y de manera explícita la probabilidad $\alpha$ de cometer ese error particular, fijando su valor antes de ver los datos. La probabilidad $\beta$, en cambio, se analiza indirectamente a través del concepto de potencia, que desarrollamos a continuación.
\end{observacion}

\subsection{Potencia de una prueba}

\begin{definicion}
La \textbf{potencia} de una prueba de hipótesis, para un valor particular $\theta_1\neq\theta_0$ del parámetro dentro de $H_1$, es la probabilidad de rechazar $H_0$ cuando el verdadero valor del parámetro es $\theta_1$:
\[
\text{Potencia}(\theta_1) = 1-\beta(\theta_1) = P(\text{Rechazar } H_0 \mid \theta=\theta_1).
\]
\end{definicion}

La potencia mide, entonces, la capacidad de la prueba para \emph{detectar} que $H_0$ es falsa, cuando el verdadero parámetro se aparta de $\theta_0$ en una magnitud dada. Como función de $\theta_1$ (manteniendo fijos $n$ y $\alpha$), a la función $\theta_1\mapsto \text{Potencia}(\theta_1)$ se la llama \textbf{función de potencia} de la prueba.

Ilustremos el cálculo de la potencia en el caso más simple: la prueba $Z$ para una media con varianza conocida, $H_0:\mu=\mu_0$ vs. $H_1:\mu>\mu_0$, con región crítica $\overline{X}>\mu_0+z_\alpha\,\sigma/\sqrt n$ (que es equivalente a $Z>z_\alpha$). Si el verdadero valor de la media es $\mu_1>\mu_0$, entonces $\overline{X}\sim N(\mu_1,\sigma^2/n)$, y
\[
\text{Potencia}(\mu_1) = P\!\left(\overline{X}>\mu_0+z_\alpha\frac{\sigma}{\sqrt n}\ \middle|\ \mu=\mu_1\right)
= P\!\left(\frac{\overline{X}-\mu_1}{\sigma/\sqrt n} > z_\alpha-\frac{\mu_1-\mu_0}{\sigma/\sqrt n}\right)
= P\!\left(Z > z_\alpha - \frac{(\mu_1-\mu_0)\sqrt n}{\sigma}\right),
\]
donde ahora $Z\sim N(0,1)$. Esta expresión permite leer directamente los tres factores que determinan la potencia:

\begin{itemize}[nosep]
\item Cuanto mayor es la diferencia $|\mu_1-\mu_0|$ (el \textbf{efecto} a detectar) en relación con $\sigma$, más negativo se vuelve el argumento de $P(Z>\cdot)$ y mayor es la potencia: es "más fácil" detectar diferencias grandes que diferencias pequeñas.
\item A mayor tamaño de muestra $n$, el término $(\mu_1-\mu_0)\sqrt n/\sigma$ crece, y por lo tanto la potencia aumenta para una misma diferencia real $\mu_1-\mu_0$. Este es el fundamento estadístico de los cálculos de \textbf{tamaño de muestra} necesarios para garantizar una potencia mínima deseada frente a un efecto mínimo de interés.
\item A mayor $\alpha$ (una región de rechazo más amplia, con $z_\alpha$ más chico), mayor es la potencia, pero a costa de un mayor riesgo de error de tipo I. Esta es, precisamente, la manifestación de la relación inversa entre $\alpha$ y $\beta$ discutida antes.
\end{itemize}

\begin{observacion}
Entre dos pruebas construidas con el mismo nivel $\alpha$, se prefiere aquella que tiene mayor potencia (es decir, menor $\beta$) para detectar el efecto de interés. Este criterio de comparación entre pruebas —maximizar la potencia sujeto a una cota sobre $\alpha$— es la base de la teoría de pruebas óptimas (uniformemente más potentes), que excede el alcance de este curso, pero que explica por qué los estadísticos de prueba que estudiaremos (basados en $\overline X$, $S^2$, etc.) no son elecciones arbitrarias sino, en general, las mejores posibles dentro de su clase.
\end{observacion}

\begin{ejemplo}
Continuando el ejemplo anterior, supongamos $\sigma=10$, $n=36$, $\alpha=0.05$ (por lo que $z_\alpha=z_{0.05}=1.645$), y que interesa la potencia para detectar un verdadero $\mu_1=\mu_0+3$. Entonces
\[
\text{Potencia}= P\!\left(Z> 1.645 - \frac{3\sqrt{36}}{10}\right) = P(Z>1.645-1.8)=P(Z>-0.155)\approx 0.5616.
\]
Es decir, con este diseño hay solo un $56\%$ de probabilidad de detectar un verdadero aumento de $3$ unidades en la media; para lograr una potencia mayor sería necesario aumentar $n$.
\end{ejemplo}

\subsection{Estadístico de prueba y región crítica}

\begin{definicion}
El \textbf{estadístico de prueba} (o estadístico de contraste) es una función de la muestra aleatoria, $T=T(X_1,\dots,X_n)$, cuya distribución de probabilidad \textbf{bajo el supuesto de que $H_0$ es verdadera} es completamente conocida.
\end{definicion}

El estadístico de prueba resume, en un único número, la información muestral relevante para el parámetro en cuestión, y permite comparar lo efectivamente observado con lo que cabría esperar si $H_0$ fuese cierta. Valores del estadístico que resultan "poco compatibles" con esa distribución de referencia constituyen evidencia en contra de $H_0$ y conducen a su rechazo.

\begin{ejemplo}
Para probar $H_0:\mu=\mu_0$ con varianza poblacional conocida $\sigma^2$, a partir de una muestra $X_1,\dots,X_n$ de una población $N(\mu,\sigma^2)$, el estadístico de prueba natural es
\[
Z=\frac{\overline X-\mu_0}{\sigma/\sqrt n}\ \overset{H_0}{\sim}\ N(0,1),
\]
lo cual se sigue directamente de que, bajo $H_0$, $\overline X\sim N(\mu_0,\sigma^2/n)$ y por lo tanto su estandarización es $N(0,1)$.
\end{ejemplo}

\begin{definicion}
El espacio de todos los valores posibles del estadístico de prueba se particiona en dos subconjuntos complementarios:
\begin{itemize}[nosep]
\item la \textbf{región crítica} o \textbf{región de rechazo}, $RR$: el conjunto de valores del estadístico que llevan a rechazar $H_0$;
\item la \textbf{región de aceptación} (o de no rechazo), $RA=RR^c$: el conjunto de valores compatibles con $H_0$.
\end{itemize}
\end{definicion}

La región crítica se determina de modo que, \emph{bajo el supuesto de que $H_0$ es verdadera}, la probabilidad de que el estadístico caiga en $RR$ sea exactamente igual al nivel de significación $\alpha$ fijado de antemano:
\[
P(T\in RR\mid H_0) = \alpha.
\]
La forma geométrica de $RR$ (si ocupa una o dos colas de la distribución de referencia) queda determinada por la forma de $H_1$:

\begin{itemize}[nosep]
\item $H_1:\theta\neq\theta_0$ (bilateral) $\Rightarrow$ $RR$ se reparte en \textbf{ambas colas} de la distribución, generalmente con $\alpha/2$ de probabilidad en cada extremo.
\item $H_1:\theta>\theta_0$ $\Rightarrow$ $RR$ ocupa la \textbf{cola derecha}, con probabilidad $\alpha$.
\item $H_1:\theta<\theta_0$ $\Rightarrow$ $RR$ ocupa la \textbf{cola izquierda}, con probabilidad $\alpha$.
\end{itemize}

Esta correspondencia tiene una justificación intuitiva: la región crítica debe ubicarse en la zona de valores del estadístico que son "más extremos" en el sentido que indica $H_1$, es decir, en la dirección en la que un valor observado sería más difícil de explicar si $H_0$ fuera cierta y más consistente con $H_1$.

A continuación se ilustra la diferencia entre una región crítica bilateral y una unilateral para el caso $Z\sim N(0,1)$ con $\alpha=0.05$.

\begin{center}
\begin{tikzpicture}[scale=0.9]
\begin{axis}[
  width=7cm, height=4cm,
  axis lines=left, ytick=\empty,
  xtick={-1.96,0,1.96}, xticklabels={$-z_{\alpha/2}$,$0$,$z_{\alpha/2}$},
  title={\footnotesize Bilateral: $H_1:\mu\neq\mu_0$},
  ymin=0, ymax=0.45,
]
\addplot[domain=-3.5:3.5, samples=100, thick, blue] {exp(-x^2/2)/sqrt(2*pi)};
\addplot[domain=-3.5:-1.96, samples=40, fill=red!40, draw=none] {exp(-x^2/2)/sqrt(2*pi)} \closedcycle;
\addplot[domain=1.96:3.5, samples=40, fill=red!40, draw=none] {exp(-x^2/2)/sqrt(2*pi)} \closedcycle;
\end{axis}
\end{tikzpicture}
\hspace{0.4cm}
\begin{tikzpicture}[scale=0.9]
\begin{axis}[
  width=7cm, height=4cm,
  axis lines=left, ytick=\empty,
  xtick={0,1.645}, xticklabels={$0$,$z_{\alpha}$},
  title={\footnotesize Unilateral: $H_1:\mu>\mu_0$},
  ymin=0, ymax=0.45,
]
\addplot[domain=-3.5:3.5, samples=100, thick, blue] {exp(-x^2/2)/sqrt(2*pi)};
\addplot[domain=1.645:3.5, samples=40, fill=red!40, draw=none] {exp(-x^2/2)/sqrt(2*pi)} \closedcycle;
\end{axis}
\end{tikzpicture}
\end{center}

En ambos gráficos, las regiones sombreadas en rojo son las regiones críticas, cuya área total (bajo la curva de densidad de $H_0$) es igual a $\alpha$ en cada caso.

\subsection{El valor p}

Además de comparar el estadístico observado con los valores críticos, existe una forma equivalente —y más informativa— de resumir la evidencia muestral en contra de $H_0$: el valor p.

\begin{definicion}
El \textbf{valor p} (o $p$-valor) es la probabilidad, calculada bajo el supuesto de que $H_0$ es verdadera, de obtener un valor del estadístico de prueba tan extremo o más extremo (en el sentido indicado por $H_1$) que el efectivamente observado en la muestra.
\end{definicion}

Formalmente, si $t_{obs}$ es el valor observado del estadístico $T$, el valor p se calcula como:
\begin{itemize}[nosep]
\item $H_1:\theta>\theta_0$: \quad $p=P(T\ge t_{obs}\mid H_0)$;
\item $H_1:\theta<\theta_0$: \quad $p=P(T\le t_{obs}\mid H_0)$;
\item $H_1:\theta\neq\theta_0$: \quad $p=2\min\{P(T\ge t_{obs}\mid H_0),\,P(T\le t_{obs}\mid H_0)\}$, que en el caso de estadísticos con distribución simétrica bajo $H_0$ (como $Z$ o $T$) se reduce a $p=2P(T\ge|t_{obs}|\mid H_0)$.
\end{itemize}

\begin{propiedad}
La regla de decisión basada en el valor p es equivalente a la regla de decisión basada en la región crítica:
\[
p\text{-valor}\le\alpha \ \Longleftrightarrow \ \text{el estadístico observado cae en la región crítica de nivel } \alpha.
\]
Por lo tanto,
\[
p\text{-valor}\le \alpha \ \Rightarrow \ \text{se rechaza } H_0; \qquad\qquad p\text{-valor}>\alpha \ \Rightarrow \ \text{no se rechaza } H_0.
\]
\end{propiedad}

\begin{proof}
Consideremos el caso unilateral derecho, $H_1:\theta>\theta_0$, con región crítica $RR=\{T>c_\alpha\}$ donde $c_\alpha$ es el valor crítico que satisface $P(T>c_\alpha\mid H_0)=\alpha$. Por definición, $p=P(T\ge t_{obs}\mid H_0)$. Como la función $c\mapsto P(T> c\mid H_0)$ es no creciente en $c$ (es la función de supervivencia de $T$ bajo $H_0$), se tiene que $t_{obs}>c_\alpha$ si y solo si $P(T>t_{obs}\mid H_0) < P(T>c_\alpha\mid H_0)=\alpha$, es decir, $t_{obs}$ cae en la región crítica si y solo si $p\le\alpha$ (salvo por la distinción entre desigualdad estricta y no estricta, irrelevante para distribuciones continuas). El mismo argumento, con las adaptaciones obvias, se aplica a los casos unilateral izquierdo y bilateral.
\end{proof}

Esta equivalencia es la razón por la cual, en la práctica —y especialmente al usar software estadístico—, suele reportarse directamente el valor p en lugar de (o además de) los valores críticos: alcanza con comparar un único número contra $\alpha$.

\begin{observacion}[Interpretación correcta del valor p]
El valor p es, posiblemente, el concepto más frecuentemente malinterpretado de la estadística inferencial. Conviene ser precisos:
\begin{itemize}[nosep]
\item El valor p \textbf{no} es la probabilidad de que $H_0$ sea verdadera, ni $1-p$ es la probabilidad de que $H_1$ sea verdadera. El valor p se calcula \emph{asumiendo que $H_0$ es cierta}; no puede, por lo tanto, ser una afirmación de probabilidad sobre la veracidad de $H_0$ misma (eso requeriría un enfoque bayesiano, con una distribución de probabilidad a priori sobre $\theta$).
\item El valor p \textbf{no} mide el tamaño o la importancia práctica del efecto detectado. Con una muestra suficientemente grande, incluso una diferencia mínima y sin relevancia práctica puede producir un valor p muy pequeño.
\item Lo que el valor p sí mide es cuán compatibles son los datos observados con $H_0$: cuanto menor es el valor p, menos compatibles son los datos con $H_0$, y por lo tanto mayor es la evidencia en contra de $H_0$ que aportan los datos.
\end{itemize}
\end{observacion}

\subsection{Procedimiento general de una prueba de hipótesis}

Formalizamos ahora el procedimiento general, que se aplicará de manera sistemática a cada una de las pruebas específicas estudiadas en las secciones siguientes. Consta de cinco pasos esenciales (que a veces se desglosan en más para mayor claridad expositiva):

\begin{enumerate}[label=\textbf{Paso \arabic*.}, wide]
\item \textbf{Formular las hipótesis.} Se plantean $H_0$ y $H_1$ en términos del parámetro poblacional de interés, de acuerdo con la pregunta de investigación, \emph{antes} de observar los datos muestrales.
\item \textbf{Fijar el nivel de significación $\alpha$ y elegir el estadístico de prueba.} Se determina el nivel de riesgo tolerado para el error de tipo I, y se elige el estadístico $T$ adecuado a la situación (parámetro de interés, supuestos sobre la población, tamaño de muestra), estableciendo su distribución exacta o aproximada bajo $H_0$.
\item \textbf{Determinar la región crítica.} De acuerdo con $\alpha$ y con la forma de $H_1$ (bilateral, unilateral derecha o unilateral izquierda), se determina el o los valores críticos que delimitan $RR$; alternativamente (o en forma complementaria), se deja planteado el cálculo del valor p.
\item \textbf{Calcular el estadístico a partir de los datos.} Se evalúa $T$ con los valores efectivamente observados en la muestra, obteniendo $t_{obs}$; y, si corresponde, se calcula el valor p asociado.
\item \textbf{Decidir y concluir.} Se rechaza $H_0$ si $t_{obs}\in RR$ (equivalentemente, si el valor p es menor o igual que $\alpha$); en caso contrario, no se rechaza $H_0$. Finalmente, la decisión estadística se traduce a una \textbf{conclusión} en el lenguaje del problema original, sin jerga técnica.
\end{enumerate}

\begin{observacion}
Algunos puntos merecen atención especial:
\begin{itemize}[nosep]
\item ``No rechazar $H_0$'' \textbf{no equivale} a ``aceptar $H_0$ como verdadera''. Significa, únicamente, que la evidencia muestral disponible no alcanza para rechazarla al nivel $\alpha$ elegido; podría deberse tanto a que $H_0$ es efectivamente cierta como a que la muestra fue insuficiente (poca potencia) para detectar un apartamiento real.
\item El nivel $\alpha$, así como la forma de $H_1$ y el tamaño de muestra $n$, deben fijarse \textbf{antes} de observar los datos. Modificar cualquiera de estos elementos \emph{después} de ver los resultados (una práctica a veces llamada informalmente \emph{$p$-hacking}) invalida las propiedades probabilísticas de la prueba, porque la probabilidad de error de tipo I real deja de ser la nominal $\alpha$.
\item El esquema conceptual (hipótesis, estadístico, distribución bajo $H_0$, región crítica, valor p, decisión) es común a \emph{todas} las pruebas específicas que se desarrollan en las secciones siguientes; lo único que cambia de una prueba a otra es el estadístico particular utilizado y su distribución de referencia.
\end{itemize}
\end{observacion}

\subsection{Dualidad entre prueba de hipótesis e intervalo de confianza}
\label{sec:dualidad}

Existe una relación profunda —y muy útil en la práctica— entre las pruebas de hipótesis bilaterales y los intervalos de confianza construidos con el mismo estadístico pivotal.

\begin{teorema}[Dualidad prueba--intervalo]
Sea $\theta$ un parámetro poblacional y sea $IC_{1-\alpha}(\theta)$ un intervalo de confianza de nivel $1-\alpha$ para $\theta$, construido a partir de un estadístico pivotal cuya distribución es simétrica y conocida. Entonces, para la prueba bilateral $H_0:\theta=\theta_0$ vs. $H_1:\theta\neq\theta_0$ al nivel de significación $\alpha$, se cumple
\[
\text{Se rechaza } H_0 \text{ al nivel } \alpha \quad \Longleftrightarrow \quad \theta_0\notin IC_{1-\alpha}(\theta).
\]
\end{teorema}

\begin{proof}
Demostramos la equivalencia en el caso de la media con varianza $\sigma^2$ conocida; el argumento se traslada sin cambios estructurales a los demás parámetros. El intervalo de confianza de nivel $1-\alpha$ para $\mu$ es
\[
IC_{1-\alpha}(\mu) = \left(\overline X - z_{\alpha/2}\frac{\sigma}{\sqrt n},\ \ \overline X + z_{\alpha/2}\frac{\sigma}{\sqrt n}\right),
\]
que se obtiene a partir del pivote $Z=\dfrac{\overline X-\mu}{\sigma/\sqrt n}\sim N(0,1)$ despejando $\mu$ en la desigualdad $-z_{\alpha/2}<Z<z_{\alpha/2}$.

Por otro lado, la región crítica de la prueba bilateral $H_0:\mu=\mu_0$ vs. $H_1:\mu\neq\mu_0$ al nivel $\alpha$, con estadístico $Z_0=\dfrac{\overline X-\mu_0}{\sigma/\sqrt n}$, es $RR=\{|Z_0|>z_{\alpha/2}\}$. Entonces:
\[
\mu_0\notin IC_{1-\alpha}(\mu)
\iff \mu_0 < \overline X - z_{\alpha/2}\frac{\sigma}{\sqrt n} \ \text{ o } \ \mu_0>\overline X+z_{\alpha/2}\frac{\sigma}{\sqrt n}
\]
\[
\iff \overline X-\mu_0 > z_{\alpha/2}\frac{\sigma}{\sqrt n} \ \text{ o } \ \overline X - \mu_0 < -z_{\alpha/2}\frac{\sigma}{\sqrt n}
\iff \left|\frac{\overline X-\mu_0}{\sigma/\sqrt n}\right| > z_{\alpha/2}
\iff |Z_0|>z_{\alpha/2},
\]
que es exactamente la condición de rechazo de $H_0$. Esto prueba la equivalencia.
\end{proof}

\begin{ejemplo}
Si un intervalo de confianza del $95\%$ para $\mu$ resulta $IC=(48.3;\,51.7)$ y se desea probar $H_0:\mu=50$ vs. $H_1:\mu\neq50$ con $\alpha=0.05$: como $50\in(48.3;51.7)$, no se rechaza $H_0$. En cambio, si se quisiera probar $H_0:\mu=52$ vs. $H_1:\mu\neq52$ con el mismo $\alpha$, como $52\notin(48.3;51.7)$, se rechazaría $H_0$.
\end{ejemplo}

\begin{observacion}
La equivalencia demostrada es válida, tal como está enunciada, para pruebas \textbf{bilaterales} y su correspondiente intervalo de confianza (bilateral) de nivel $1-\alpha$. Para pruebas \textbf{unilaterales}, existe una correspondencia análoga con intervalos de confianza de una cola (semirrectas): la prueba $H_0:\mu=\mu_0$ vs. $H_1:\mu>\mu_0$ al nivel $\alpha$ es equivalente a rechazar $H_0$ cuando $\mu_0$ queda por debajo del extremo inferior de un intervalo de confianza unilateral de nivel $1-\alpha$ de la forma $\left(\overline X - z_\alpha \sigma/\sqrt n,\ +\infty\right)$.

La utilidad práctica de esta dualidad es considerable: un único intervalo de confianza permite decidir simultáneamente sobre \emph{cualquier} valor hipotético $\theta_0$ (bilateral), mientras que cada prueba de hipótesis, tal como se plantea habitualmente, evalúa un único valor $\theta_0$ a la vez. Por esta razón, cuando se dispone de un intervalo de confianza ya calculado, suele ser más económico usarlo directamente para responder preguntas de prueba de hipótesis que recalcular el estadístico de prueba desde cero.
\end{observacion}

\newpage
\section{Pruebas para una media y una varianza poblacional}

En esta sección aplicamos el marco general desarrollado a los dos parámetros más habituales de una única población: la media $\mu$ y la varianza $\sigma^2$. Como se anticipó, lo único que cambia de una prueba a otra es el estadístico utilizado, su distribución bajo $H_0$ y, en consecuencia, la forma de la región crítica; la lógica de decisión (pasos 1, 2, 5) permanece invariable.

\subsection{Prueba para la media con varianza conocida (prueba Z)}

Sea $X_1,\dots,X_n$ una muestra aleatoria de una población $N(\mu,\sigma^2)$ (o, más generalmente, de cualquier población con varianza $\sigma^2$ finita y $n$ suficientemente grande, apelando al Teorema Central del Límite), con $\sigma^2$ \textbf{conocida}. Interesa probar
\[
H_0:\mu=\mu_0.
\]

\begin{teorema}[Estadístico de prueba]
Bajo $H_0$, el estadístico
\[
Z=\frac{\overline X-\mu_0}{\sigma/\sqrt n}
\]
se distribuye exactamente $N(0,1)$ si la población es normal, y aproximadamente $N(0,1)$ para $n$ grande en el caso general, por el Teorema Central del Límite.
\end{teorema}

\begin{proof}
Si $X_1,\dots,X_n$ son i.i.d. $N(\mu,\sigma^2)$, entonces $\overline X$ es una combinación lineal de variables normales independientes, por lo que $\overline X\sim N\!\left(\mu,\ \sigma^2/n\right)$ exactamente, para cualquier $n$. Bajo $H_0$, $\mu=\mu_0$, luego $\overline X\sim N(\mu_0,\sigma^2/n)$, y estandarizando, $Z=(\overline X-\mu_0)/(\sigma/\sqrt n)\sim N(0,1)$. Cuando la población no es normal pero $n$ es grande, el Teorema Central del Límite garantiza que $\overline X$ es aproximadamente $N(\mu,\sigma^2/n)$, y el mismo argumento produce la distribución aproximada de $Z$.
\end{proof}

La región crítica se construye de modo que, bajo $H_0$, el estadístico $Z$ caiga en ella con probabilidad exactamente $\alpha$, respetando el sentido de $H_1$:

\begin{table}[h]
\centering
\small
\begin{tabular}{@{}lcc@{}}
\toprule
$H_1$ & Región crítica & Naturaleza \\
\midrule
$\mu\neq\mu_0$ & $|Z|>z_{\alpha/2}$ & bilateral \\
$\mu>\mu_0$ & $Z>z_{\alpha}$ & cola derecha \\
$\mu<\mu_0$ & $Z<-z_{\alpha}$ & cola izquierda \\
\bottomrule
\end{tabular}
\end{table}

\noindent donde $z_\gamma$ denota el valor tal que $P(Z>z_\gamma)=\gamma$ para $Z\sim N(0,1)$. Estas regiones se deducen exigiendo que la probabilidad de rechazo bajo $H_0$ sea $\alpha$: por ejemplo, en el caso bilateral, $P(|Z|>z_{\alpha/2}\mid H_0)=P(Z>z_{\alpha/2})+P(Z<-z_{\alpha/2})=\alpha/2+\alpha/2=\alpha$, usando la simetría de la $N(0,1)$.

\begin{ejemplo}[Prueba bilateral, deducción completa]
Una máquina envasa paquetes de café cuyo peso se distribuye $N(\mu,\sigma^2)$ con $\sigma=10$ g. El rótulo indica $\mu_0=500$ g. Se toma una muestra de $n=36$ paquetes y se obtiene $\overline x=496$ g. Con $\alpha=0.05$, ¿hay evidencia de que el peso medio difiere de $500$ g?

\textbf{1. Hipótesis:} $H_0:\mu=500$ vs. $H_1:\mu\neq500$.

\textbf{2. Estadístico:} $\displaystyle Z=\frac{\overline X-500}{10/\sqrt{36}}\overset{H_0}{\sim}N(0,1)$.

\textbf{3. Región crítica} ($\alpha=0.05$, bilateral): $|Z|>z_{0.025}=1.96$.

\textbf{4. Cálculo:}
\[
z=\frac{496-500}{10/\sqrt{36}}=\frac{-4}{1.6667}=-2.40.
\]

\textbf{5. Decisión por región crítica:} como $|-2.40|=2.40>1.96$, se \textbf{rechaza} $H_0$.

\textbf{6. Cálculo del valor p (verificación):} para la prueba bilateral,
\[
p=2\,P(Z<-2.40)=2\times0.0082=0.0164.
\]
Como $p=0.0164<\alpha=0.05$, se llega exactamente a la misma decisión: se rechaza $H_0$, confirmando la equivalencia entre ambos criterios demostrada en la sección anterior.

\textbf{7. Conclusión:} con un $5\%$ de significación, hay evidencia estadísticamente significativa de que el peso medio de los paquetes difiere de $500$ g.
\end{ejemplo}

\begin{observacion}
Si en el ejemplo anterior la hipótesis alternativa hubiera sido unilateral, $H_1:\mu<500$, la región crítica sería $Z<-z_{0.05}=-1.645$; como $-2.40<-1.645$, también se rechazaría $H_0$ (con más fuerza, dado que la región crítica unilateral es "más generosa" en esa cola), y el valor p sería $p=P(Z<-2.40)=0.0082$, la mitad del valor p bilateral. Esta relación —el valor p unilateral es la mitad del bilateral, cuando el signo observado coincide con el sentido de $H_1$— es general para estadísticos con distribución simétrica.
\end{observacion}

\subsection{Prueba para la media con varianza desconocida (prueba t)}

En la inmensa mayoría de las aplicaciones reales, $\sigma^2$ no se conoce y debe estimarse a partir de la propia muestra mediante la cuasivarianza muestral
\[
S^2=\frac{1}{n-1}\sum_{i=1}^n (X_i-\overline X)^2.
\]
Al reemplazar $\sigma$ por $S$ en el estadístico $Z$, se introduce variabilidad adicional (la propia de estimar $\sigma$), que hace que la distribución exacta ya no sea normal estándar.

\begin{teorema}
Sea $X_1,\dots,X_n$ una muestra aleatoria de una población $N(\mu,\sigma^2)$. Bajo $H_0:\mu=\mu_0$, el estadístico
\[
T=\frac{\overline X-\mu_0}{S/\sqrt n}
\]
se distribuye exactamente según una $t$ de Student con $n-1$ grados de libertad, $t_{n-1}$.
\end{teorema}

\begin{proof}[Idea de la demostración]
Se sabe, de la Unidad 5, que para una muestra normal, $\overline X$ y $S^2$ son independientes, que $\overline X\sim N(\mu,\sigma^2/n)$ y que $(n-1)S^2/\sigma^2\sim\chi^2_{n-1}$. Por definición, si $Z\sim N(0,1)$ y $W\sim\chi^2_{k}$ son independientes, entonces $Z/\sqrt{W/k}\sim t_k$. Tomando $Z=(\overline X-\mu_0)/(\sigma/\sqrt n)$ (que bajo $H_0$ es $N(0,1)$) y $W=(n-1)S^2/\sigma^2\sim\chi^2_{n-1}$, ambos independientes, se obtiene
\[
\frac{Z}{\sqrt{W/(n-1)}} = \frac{\dfrac{\overline X-\mu_0}{\sigma/\sqrt n}}{\sqrt{\dfrac{(n-1)S^2/\sigma^2}{n-1}}} = \frac{\dfrac{\overline X-\mu_0}{\sigma/\sqrt n}}{S/\sigma} = \frac{\overline X-\mu_0}{S/\sqrt n} = T,
\]
donde el factor $\sigma$ se cancela algebraicamente. Por lo tanto $T\sim t_{n-1}$ bajo $H_0$.
\end{proof}

\begin{table}[h]
\centering
\small
\begin{tabular}{@{}lc@{}}
\toprule
$H_1$ & Región crítica \\
\midrule
$\mu\neq\mu_0$ & $|T|>t_{n-1,\,\alpha/2}$ \\
$\mu>\mu_0$ & $T>t_{n-1,\,\alpha}$ \\
$\mu<\mu_0$ & $T<-t_{n-1,\,\alpha}$ \\
\bottomrule
\end{tabular}
\end{table}

\begin{observacion}
Cuando $n$ es grande ($n\ge30$, criterio usual), la distribución $t_{n-1}$ se aproxima cada vez más a la $N(0,1)$ (los cuantiles $t_{n-1,\gamma}\to z_\gamma$ cuando $n\to\infty$), y ambos procedimientos son prácticamente equivalentes en la práctica. Sin embargo, conceptualmente, si $\sigma^2$ no se conoce, el estadístico correcto sigue siendo $T$ con $n-1$ grados de libertad, y no $Z$; usar $Z$ subestima levemente la incertidumbre real cuando $n$ no es muy grande.
\end{observacion}

\begin{ejemplo}[Prueba unilateral con valor p]
Un fabricante de neumáticos afirma que la duración media de un nuevo modelo supera las $40\,000$ km. En una muestra de $n=16$ neumáticos (duración $\sim N(\mu,\sigma^2)$) se obtuvo $\overline x=41\,200$ km y $s=1500$ km. Con $\alpha=0.05$, ¿respaldan los datos la afirmación del fabricante?

\textbf{1. Hipótesis:} $H_0:\mu\le40\,000$ vs. $H_1:\mu>40\,000$ (se trabaja con $H_0:\mu=40\,000$, el punto frontera).

\textbf{2. Estadístico:} $\displaystyle T=\frac{\overline X-40\,000}{S/\sqrt{16}}\overset{H_0}{\sim} t_{15}$.

\textbf{3. Región crítica:} $T>t_{15,\,0.05}=1.753$.

\textbf{4. Cálculo:}
\[
t=\frac{41\,200-40\,000}{1500/\sqrt{16}}=\frac{1200}{375}=3.20.
\]

\textbf{5. Decisión por región crítica:} $3.20>1.753\Rightarrow$ se rechaza $H_0$.

\textbf{6. Valor p:} usando tablas de $t_{15}$, $P(T_{15}>3.20)$ se ubica entre $0.0025$ y $0.005$ (los valores tabulados usuales dan $t_{15,0.005}=2.947$ y $t_{15,0.0025}\approx3.286$); concretamente, $p\approx0.0029$. Como $p<0.05$, coincide con la decisión anterior.

\textbf{7. Conclusión:} con un $5\%$ de significación, la evidencia muestral respalda que la duración media supera los $40\,000$ km.
\end{ejemplo}

\begin{ejemplo}[Prueba bilateral con no rechazo]
Un laboratorio farmacéutico debe verificar que los comprimidos contienen $\mu_0=250$ mg del principio activo. En $n=10$ comprimidos se midió $\overline x=247.5$ mg y $s=4.8$ mg. Con $\alpha=0.05$:

$H_0:\mu=250$ vs. $H_1:\mu\neq250$. Región crítica (gl$=9$): $|T|>t_{9,0.025}=2.262$.
\[
t=\frac{247.5-250}{4.8/\sqrt{10}}=\frac{-2.5}{1.518}=-1.65.
\]
Como $|-1.65|=1.65<2.262$, no se rechaza $H_0$: la evidencia muestral no es suficiente para afirmar que el contenido medio difiere de $250$ mg. El valor p bilateral, $p=2P(T_9<-1.65)\approx 2\times0.0665= 0.133$, confirma que la evidencia en contra de $H_0$ es débil (bastante mayor que $\alpha=0.05$).
\end{ejemplo}

\subsection{Prueba para la varianza poblacional}

Sea nuevamente $X_1,\dots,X_n$ una muestra aleatoria de $N(\mu,\sigma^2)$. Interesa ahora probar afirmaciones sobre la variabilidad de la población, no sobre su valor central:
\[
H_0:\sigma^2=\sigma_0^2.
\]

\begin{teorema}[Estadístico de prueba para $\sigma^2$]
Bajo $H_0$, el estadístico
\[
\chi^2=\frac{(n-1)S^2}{\sigma_0^2}
\]
se distribuye exactamente $\chi^2_{n-1}$.
\end{teorema}

\begin{proof}
Es un resultado establecido en la Unidad 5 (distribuciones asociadas al muestreo normal) que, si $X_1,\dots,X_n$ es una muestra aleatoria de $N(\mu,\sigma^2)$, entonces
\[
\frac{(n-1)S^2}{\sigma^2}\sim\chi^2_{n-1},
\]
lo cual se deduce de escribir $(n-1)S^2/\sigma^2 = \sum_{i=1}^n\left(\frac{X_i-\overline X}{\sigma}\right)^2$ y observar que esta es una forma cuadrática en variables normales estándar con $n-1$ grados de libertad efectivos, debido a la restricción lineal $\sum_i (X_i-\overline X)=0$ que reduce en uno la dimensión del espacio de variación libre (formalmente, se sigue del teorema de Cochran o de una diagonalización ortogonal del vector $(X_i-\overline X)/\sigma$). Bajo $H_0:\sigma^2=\sigma_0^2$, sustituyendo $\sigma^2$ por $\sigma_0^2$ se obtiene el estadístico enunciado, que hereda la distribución $\chi^2_{n-1}$.
\end{proof}

A diferencia de $Z$ y $T$, la distribución $\chi^2_{n-1}$ \textbf{no es simétrica}: es asimétrica a derecha y toma únicamente valores no negativos. Esto tiene una consecuencia importante para el caso bilateral: los dos valores críticos que delimitan la región de rechazo \emph{no} son opuestos entre sí (no hay noción de "signo" en $\chi^2$), sino que se obtienen por separado de las colas inferior y superior de la distribución.

\begin{table}[h]
\centering
\small
\begin{tabular}{@{}lc@{}}
\toprule
$H_1$ & Región crítica \\
\midrule
$\sigma^2\neq\sigma_0^2$ & $\chi^2<\chi^2_{n-1,\,1-\alpha/2}$ \ o \ $\chi^2>\chi^2_{n-1,\,\alpha/2}$ \\
$\sigma^2>\sigma_0^2$ & $\chi^2>\chi^2_{n-1,\,\alpha}$ \\
$\sigma^2<\sigma_0^2$ & $\chi^2<\chi^2_{n-1,\,1-\alpha}$ \\
\bottomrule
\end{tabular}
\end{table}

\noindent donde $\chi^2_{n-1,\gamma}$ es el valor tal que $P(\chi^2_{n-1}>\chi^2_{n-1,\gamma})=\gamma$.

\begin{center}
\begin{tikzpicture}
\begin{axis}[
  width=10cm, height=5cm,
  axis lines=left, ytick=\empty,
  xlabel={$\chi^2$}, ylabel={densidad},
  xmin=0, xmax=40, ymin=0,
  legend pos=north east, legend style={font=\footnotesize},
]
\addplot[domain=0.1:40, samples=150, thick, blue] {x^(2)*exp(-x/2)};
\addlegendentry{$\chi^2_{5}$ (asimétrica, cola derecha larga)}
\addplot[domain=0.1:40, samples=150, thick, red, dashed] {x^(7)*exp(-x/2)*0.02};
\addlegendentry{$\chi^2_{15}$ (más simétrica al crecer gl)}
\end{axis}
\end{tikzpicture}
\end{center}

\begin{ejemplo}[Prueba unilateral derecha, con valor p]
Un proceso de llenado de botellas debe tener una desviación estándar no mayor a $\sigma_0=2$ mL (varianza $\sigma_0^2=4$). En una muestra de $n=20$ botellas se midió $s^2=6.1$ mL$^2$. Con $\alpha=0.05$, ¿hay evidencia de que la variabilidad del proceso excede la tolerancia?

\textbf{1. Hipótesis:} $H_0:\sigma^2\le4$ vs. $H_1:\sigma^2>4$ (se usa $H_0:\sigma^2=4$).

\textbf{2. Estadístico:} $\displaystyle \chi^2=\frac{(n-1)S^2}{4}\overset{H_0}{\sim}\chi^2_{19}$.

\textbf{3. Región crítica:} $\chi^2>\chi^2_{19,\,0.05}=30.14$.

\textbf{4. Cálculo:}
\[
\chi^2=\frac{19\times6.1}{4}=\frac{115.9}{4}=28.98.
\]

\textbf{5. Decisión:} $28.98<30.14\Rightarrow$ no se rechaza $H_0$.

\textbf{6. Valor p:} de tablas de $\chi^2_{19}$, $\chi^2_{19,0.10}=27.20$ y $\chi^2_{19,0.05}=30.14$; como $27.20<28.98<30.14$, el valor p está entre $0.05$ y $0.10$ (aproximadamente $p\approx0.067$), mayor que $\alpha=0.05$, confirmando el no rechazo.

\textbf{7. Conclusión:} con un $5\%$ de significación, no hay evidencia suficiente de que la variabilidad del proceso supere la tolerancia especificada, aunque el resultado está cerca del límite.
\end{ejemplo}

\begin{ejemplo}[Prueba bilateral]
Un laboratorio requiere que la varianza de las mediciones de un instrumento sea exactamente $\sigma_0^2=0.25$. Una muestra de $n=15$ mediciones da $s^2=0.42$. Con $\alpha=0.10$, ¿hay evidencia de que la varianza real difiere de $0.25$?

\textbf{Hipótesis:} $H_0:\sigma^2=0.25$ vs. $H_1:\sigma^2\neq0.25$.

\textbf{Región crítica} ($\alpha=0.10$, $14$ gl): se rechaza si $\chi^2<\chi^2_{14,0.95}=6.571$ o $\chi^2>\chi^2_{14,0.05}=23.68$.

\textbf{Cálculo:}
\[
\chi^2=\frac{14\times0.42}{0.25}=\frac{5.88}{0.25}=23.52.
\]

\textbf{Decisión:} como $23.52<23.68$, no se rechaza $H_0$, aunque el valor está muy cerca del límite crítico. Consistentemente, el valor p bilateral (aproximadamente $p\approx2\times0.052=0.104$, ligeramente mayor que $0.10$) confirma la decisión, aunque por un margen muy estrecho.

\textbf{Conclusión:} al nivel $10\%$, no hay evidencia suficiente para afirmar que la varianza difiere de $0.25$; el resultado sugiere revisar el proceso con una muestra mayor antes de descartar definitivamente un cambio en la variabilidad.
\end{ejemplo}

\begin{observacion}[Sensibilidad a la normalidad]
A diferencia de las pruebas $Z$ y $T$ para la media (que, por el TCL, son razonablemente robustas ante apartamientos moderados de la normalidad cuando $n$ es grande), la prueba $\chi^2$ para la varianza es \textbf{muy sensible} al supuesto de normalidad poblacional: si la población tiene curtosis distinta de la normal, la distribución exacta de $(n-1)S^2/\sigma_0^2$ se aparta considerablemente de $\chi^2_{n-1}$, incluso para $n$ grande. Por esta razón, esta prueba debe usarse con cautela cuando hay evidencia de que la población no es aproximadamente normal.
\end{observacion}

\newpage
\section{Pruebas para dos muestras}

Es habitual comparar dos poblaciones a través de sus medias o de sus varianzas: dos procesos de producción, dos proveedores, un tratamiento frente a un control, mediciones antes y después de una intervención, etc. Antes de elegir el procedimiento adecuado, es imprescindible clasificar correctamente el diseño del estudio en una de las siguientes categorías:

\begin{itemize}[nosep]
\item \textbf{Muestras independientes} con varianzas poblacionales \textbf{conocidas}.
\item \textbf{Muestras independientes} con varianzas poblacionales \textbf{desconocidas pero supuestas iguales} (se estima una varianza combinada o \emph{pooled}).
\item \textbf{Muestras apareadas} (o dependientes): cada observación de una muestra está naturalmente asociada a una observación de la otra (mismo individuo medido antes y después, mismo lote medido con dos instrumentos distintos, gemelos, etc.).
\item Comparación de las \textbf{varianzas} de dos poblaciones (prueba $F$).
\end{itemize}

\begin{observacion}
La pregunta clave del diseño, previa a cualquier cálculo, es: ¿existe una correspondencia natural entre cada observación de una muestra y una de la otra? Ejemplos de diseños \textbf{independientes}: dos máquinas distintas, dos proveedores distintos, grupo control vs. grupo tratado con sujetos distintos. Ejemplos de diseños \textbf{apareados}: mismo sujeto antes y después de un tratamiento, misma pieza medida por dos instrumentos, unidades emparejadas por algún criterio de similitud. Confundir el diseño lleva a elegir un estadístico incorrecto —típicamente, tratar como independientes datos que en realidad están apareados, lo que produce estimaciones de varianza infladas y pruebas con menor potencia de la que realmente se podría lograr— y, en consecuencia, a conclusiones erróneas.
\end{observacion}

\subsection{Diferencia de medias: varianzas conocidas}

Sean $X_1,\dots,X_{n_1}$ una muestra de $N(\mu_1,\sigma_1^2)$ y $Y_1,\dots,Y_{n_2}$ una muestra de $N(\mu_2,\sigma_2^2)$, \textbf{independientes} entre sí, con $\sigma_1^2$ y $\sigma_2^2$ \textbf{conocidas}. Se desea probar
\[
H_0:\mu_1-\mu_2=\delta_0 \qquad (\text{usualmente } \delta_0=0).
\]

\begin{teorema}
Bajo $H_0$, el estadístico
\[
Z=\frac{(\overline X-\overline Y)-\delta_0}{\sqrt{\dfrac{\sigma_1^2}{n_1}+\dfrac{\sigma_2^2}{n_2}}}
\]
se distribuye $N(0,1)$.
\end{teorema}

\begin{proof}
Por independencia de las dos muestras, $\overline X\sim N(\mu_1,\sigma_1^2/n_1)$ y $\overline Y\sim N(\mu_2,\sigma_2^2/n_2)$ son independientes entre sí, y por lo tanto su diferencia también es normal, con
\[
E[\overline X-\overline Y]=\mu_1-\mu_2, \qquad \text{Var}(\overline X-\overline Y)=\text{Var}(\overline X)+\text{Var}(\overline Y)=\frac{\sigma_1^2}{n_1}+\frac{\sigma_2^2}{n_2}
\]
(la suma de varianzas se debe a la independencia, que anula el término de covarianza). Bajo $H_0:\mu_1-\mu_2=\delta_0$, estandarizando $\overline X-\overline Y$ con esta media y varianza se obtiene exactamente el estadístico $Z\sim N(0,1)$ enunciado.
\end{proof}

Las regiones críticas son enteramente análogas al caso de una sola media: $|Z|>z_{\alpha/2}$ para $H_1:\mu_1-\mu_2\neq\delta_0$; $Z>z_\alpha$ para $H_1:\mu_1-\mu_2>\delta_0$; $Z<-z_\alpha$ para $H_1:\mu_1-\mu_2<\delta_0$.

\begin{ejemplo}
Dos máquinas envasadoras de un mismo producto tienen desviaciones estándar históricas conocidas $\sigma_1=5$ g y $\sigma_2=4$ g. Se toman $n_1=40$ envases de la máquina 1 ($\overline x=252$ g) y $n_2=35$ de la máquina 2 ($\overline y=249$ g). Con $\alpha=0.05$, ¿difiere el peso medio entregado por ambas máquinas?

\textbf{1. Hipótesis:} $H_0:\mu_1-\mu_2=0$ vs. $H_1:\mu_1-\mu_2\neq0$.

\textbf{2. Región crítica:} $|Z|>z_{0.025}=1.96$.

\textbf{3. Cálculo:}
\[
z=\frac{(252-249)-0}{\sqrt{25/40+16/35}}=\frac{3}{\sqrt{0.625+0.457}}=\frac{3}{\sqrt{1.082}}=\frac{3}{1.040}=2.88.
\]

\textbf{4. Decisión:} $2.88>1.96\Rightarrow$ se rechaza $H_0$.

\textbf{5. Valor p:} $p=2P(Z>2.88)=2\times0.0020=0.0040$, mucho menor que $\alpha=0.05$ (y también que $0.01$): la evidencia en contra de $H_0$ es muy fuerte.

\textbf{6. Intervalo de confianza asociado (verificación mediante dualidad):} un IC del $95\%$ para $\mu_1-\mu_2$ es
\[
(\overline x-\overline y)\pm z_{0.025}\sqrt{\tfrac{\sigma_1^2}{n_1}+\tfrac{\sigma_2^2}{n_2}} = 3\pm1.96\times1.040 = 3\pm2.04 = (0.96;\,5.04),
\]
que no contiene al $0$, consistente con el rechazo de $H_0:\mu_1-\mu_2=0$ de acuerdo con la dualidad prueba--intervalo demostrada en la Sección \ref{sec:dualidad}.

\textbf{7. Conclusión:} hay evidencia significativa de que el peso medio entregado por ambas máquinas difiere.
\end{ejemplo}

\subsection{Diferencia de medias: varianzas desconocidas e iguales (pooled)}

Cuando $\sigma_1^2$ y $\sigma_2^2$ son desconocidas, pero resulta razonable suponer que $\sigma_1^2=\sigma_2^2=\sigma^2$ (poblaciones normales, muestras independientes), conviene combinar la información de ambas muestras en una única estimación de la varianza común $\sigma^2$, en lugar de estimarla dos veces por separado. Esto se logra mediante la \textbf{varianza combinada} o \emph{pooled}.

\begin{definicion}[Varianza combinada]
\[
S_p^2=\frac{(n_1-1)S_1^2+(n_2-1)S_2^2}{n_1+n_2-2}.
\]
\end{definicion}

\begin{propiedad}[Justificación del estimador combinado]
$S_p^2$ es un promedio ponderado de $S_1^2$ y $S_2^2$, con pesos proporcionales a los respectivos grados de libertad $n_1-1$ y $n_2-1$, y es un estimador insesgado de $\sigma^2$ cuando $\sigma_1^2=\sigma_2^2=\sigma^2$.
\end{propiedad}

\begin{proof}
Como $(n_1-1)S_1^2/\sigma^2\sim\chi^2_{n_1-1}$ y $(n_2-1)S_2^2/\sigma^2\sim\chi^2_{n_2-1}$ son independientes (provienen de muestras independientes), su suma se distribuye $\chi^2_{n_1+n_2-2}$ (propiedad aditiva de la $\chi^2$), es decir,
\[
\frac{(n_1-1)S_1^2+(n_2-1)S_2^2}{\sigma^2}\sim\chi^2_{n_1+n_2-2}.
\]
Como la esperanza de una $\chi^2_k$ es $k$, se obtiene
\[
E\!\left[\frac{(n_1-1)S_1^2+(n_2-1)S_2^2}{\sigma^2}\right]=n_1+n_2-2 \ \Longrightarrow\ E[S_p^2]=E\!\left[\frac{(n_1-1)S_1^2+(n_2-1)S_2^2}{n_1+n_2-2}\right]=\sigma^2,
\]
lo que muestra que $S_p^2$ es insesgado para $\sigma^2$. Además, de la misma propiedad aditiva se sigue que $\dfrac{(n_1+n_2-2)S_p^2}{\sigma^2}\sim\chi^2_{n_1+n_2-2}$, resultado que se usa a continuación para derivar la distribución del estadístico de prueba.
\end{proof}

\begin{teorema}[Estadístico $t$ agrupado]
Bajo $H_0:\mu_1-\mu_2=\delta_0$, con $\sigma_1^2=\sigma_2^2$ desconocida,
\[
T=\frac{(\overline X-\overline Y)-\delta_0}{S_p\sqrt{\dfrac{1}{n_1}+\dfrac{1}{n_2}}}\ \sim\ t_{n_1+n_2-2}.
\]
\end{teorema}

\begin{proof}
Bajo $H_0$ y el supuesto $\sigma_1^2=\sigma_2^2=\sigma^2$, sabemos que $Z=\dfrac{(\overline X-\overline Y)-\delta_0}{\sigma\sqrt{1/n_1+1/n_2}}\sim N(0,1)$ (por el mismo argumento que en el caso de varianzas conocidas, con $\sigma_1^2=\sigma_2^2=\sigma^2$), y que $W=\dfrac{(n_1+n_2-2)S_p^2}{\sigma^2}\sim\chi^2_{n_1+n_2-2}$, independiente de $Z$ (pues $\overline X, \overline Y$ son independientes de $S_1^2,S_2^2$ dentro de cada muestra, y las dos muestras son independientes entre sí). Aplicando la definición de la $t$ de Student, $T=Z/\sqrt{W/(n_1+n_2-2)}$:
\[
T=\frac{\dfrac{(\overline X-\overline Y)-\delta_0}{\sigma\sqrt{1/n_1+1/n_2}}}{\sqrt{\dfrac{(n_1+n_2-2)S_p^2/\sigma^2}{n_1+n_2-2}}}=\frac{\dfrac{(\overline X-\overline Y)-\delta_0}{\sigma\sqrt{1/n_1+1/n_2}}}{S_p/\sigma}=\frac{(\overline X-\overline Y)-\delta_0}{S_p\sqrt{1/n_1+1/n_2}},
\]
donde nuevamente $\sigma$ se cancela, quedando el estadístico enunciado, distribuido $t_{n_1+n_2-2}$.
\end{proof}

Las regiones críticas son análogas a las del caso de varianzas conocidas, reemplazando $z_\gamma$ por $t_{n_1+n_2-2,\gamma}$.

\begin{ejemplo}
Se compara la resistencia a la rotura (en kg) de cables de dos proveedores, suponiendo varianzas poblacionales iguales. Proveedor A: $n_1=12$, $\overline x=436$, $s_1^2=98$. Proveedor B: $n_2=10$, $\overline y=428$, $s_2^2=110$. Con $\alpha=0.05$, ¿el proveedor A ofrece mayor resistencia media?

\textbf{1. Hipótesis:} $H_0:\mu_1-\mu_2\le0$ vs. $H_1:\mu_1-\mu_2>0$.

\textbf{2. Varianza combinada:}
\[
S_p^2=\frac{11(98)+9(110)}{20}=\frac{1078+990}{20}=\frac{2068}{20}=103.4,\qquad S_p=10.17.
\]

\textbf{3. Región crítica} (gl$=20$): $T>t_{20,0.05}=1.725$.

\textbf{4. Cálculo:}
\[
t=\frac{(436-428)-0}{10.17\sqrt{\tfrac1{12}+\tfrac1{10}}}=\frac{8}{10.17\times0.4282}=\frac{8}{4.354}=1.84.
\]

\textbf{5. Decisión:} $1.84>1.725\Rightarrow$ se rechaza $H_0$.

\textbf{6. Valor p:} de tablas de $t_{20}$, $t_{20,0.05}=1.725$ y $t_{20,0.025}=2.086$; como $1.725<1.84<2.086$, el valor p está entre $0.025$ y $0.05$ (aproximadamente $p\approx0.040$), consistente con el rechazo.

\textbf{7. Conclusión:} con un $5\%$ de significación, hay evidencia de que el proveedor A entrega cables con mayor resistencia media.
\end{ejemplo}

\begin{observacion}[Criterio para decidir entre varianzas iguales o distintas; mención de Welch]
Cuando las varianzas poblacionales son desconocidas, la validez del estadístico $t$ agrupado depende críticamente del supuesto $\sigma_1^2=\sigma_2^2$. Como criterio práctico, este supuesto puede evaluarse formalmente mediante la prueba $F$ de igualdad de varianzas que se desarrolla en la subsección siguiente; una regla informal frecuentemente usada es que si el cociente entre la mayor y la menor varianza muestral no supera aproximadamente $3$ ó $4$ (y los tamaños muestrales son similares), el estadístico agrupado suele comportarse razonablemente bien incluso ante apartamientos moderados del supuesto de igualdad.

Cuando las varianzas poblacionales son desconocidas y \textbf{no} resulta razonable suponerlas iguales, el estadístico $t$ agrupado deja de tener la distribución $t_{n_1+n_2-2}$ exacta y puede producir un nivel de significación real distinto del nominal, especialmente si los tamaños de muestra $n_1$ y $n_2$ son muy distintos. Para esa situación existe una extensión clásica, la \textbf{aproximación de Welch--Satterthwaite}, que utiliza el estadístico
\[
T'=\frac{(\overline X-\overline Y)-\delta_0}{\sqrt{S_1^2/n_1+S_2^2/n_2}},
\]
cuya distribución bajo $H_0$ se aproxima por una $t$ con grados de libertad \emph{ajustados} (no enteros, en general) dados por la fórmula de Satterthwaite
\[
\nu \approx \frac{\left(\dfrac{S_1^2}{n_1}+\dfrac{S_2^2}{n_2}\right)^{2}}{\dfrac{(S_1^2/n_1)^2}{n_1-1}+\dfrac{(S_2^2/n_2)^2}{n_2-1}}.
\]
El desarrollo completo de esta corrección excede el alcance de este curso; en esta materia se trabaja siempre bajo el supuesto de varianzas desconocidas pero \textbf{iguales} cuando corresponde usar una prueba $t$ para dos muestras independientes, contrastando previamente ese supuesto con la prueba $F$ cuando resulte pertinente.
\end{observacion}

\subsection{Muestras apareadas}

Cuando las dos muestras \textbf{no son independientes} —porque provienen de una correspondencia natural entre observaciones (mismo sujeto medido dos veces, mismo elemento evaluado por dos métodos, pares naturalmente asociados)—, no corresponde aplicar ninguno de los procedimientos anteriores. La estrategia adecuada consiste en trabajar con las \textbf{diferencias individuales}
\[
D_i=X_i-Y_i,\qquad i=1,\dots,n,
\]
reduciendo el problema de comparación de dos medias a una prueba de una sola media (la de las diferencias), aplicada sobre la Sección 2.

\begin{propiedad}
Si $D_1,\dots,D_n$ son aproximadamente $N(\mu_D,\sigma_D^2)$ con $\sigma_D^2$ desconocida, entonces bajo $H_0:\mu_D=\delta_0$ (usualmente $\delta_0=0$, es decir, $H_0:\mu_1=\mu_2$),
\[
T=\frac{\overline D-\delta_0}{S_D/\sqrt n}\ \overset{H_0}{\sim}\ t_{n-1},
\]
donde $\overline D$ y $S_D$ son, respectivamente, la media y el desvío muestral de las diferencias $D_1,\dots,D_n$.
\end{propiedad}

Esto es una consecuencia inmediata del resultado ya demostrado para la prueba $t$ de una media: basta reemplazar la muestra original por la muestra (univariada) de diferencias $D_i$.

\begin{observacion}[Por qué conviene trabajar con las diferencias]
Al trabajar con las diferencias, cada unidad experimental actúa como su propio control, y se elimina de la comparación toda la variabilidad debida a diferencias \emph{entre} unidades (por ejemplo, diferencias basales entre sujetos que nada tienen que ver con el efecto del tratamiento que se quiere detectar). Formalmente, si $X_i$ e $Y_i$ están correlacionadas positivamente (lo cual es típico cuando provienen del mismo sujeto o unidad), entonces
\[
\text{Var}(D_i)=\text{Var}(X_i-Y_i)=\sigma_1^2+\sigma_2^2-2\,\text{Cov}(X_i,Y_i) < \sigma_1^2+\sigma_2^2,
\]
que es precisamente la varianza que se usaría (erróneamente) si se tratara a las muestras como independientes ignorando el apareamiento. Es decir, ignorar una correlación positiva entre pares infla artificialmente la estimación de la varianza de la diferencia de medias, lo que reduce la potencia de la prueba. El análisis apareado, al usar $\text{Var}(D_i)$ directamente (más chica cuando la correlación es positiva), suele ser considerablemente más potente que tratar las muestras —incorrectamente— como independientes.
\end{observacion}

\begin{ejemplo}
Se evalúa un programa de entrenamiento midiendo el tiempo (en segundos) que tardan $8$ operarios en completar una tarea, antes y después del entrenamiento:

\begin{center}
\small
\begin{tabular}{@{}lcccccccc@{}}
\toprule
Operario & 1 & 2 & 3 & 4 & 5 & 6 & 7 & 8 \\
\midrule
Antes & 25 & 28 & 22 & 30 & 27 & 24 & 26 & 29 \\
Después & 23 & 24 & 22 & 27 & 25 & 24 & 22 & 25 \\
Diferencia $d_i$ & 2 & 4 & 0 & 3 & 2 & 0 & 4 & 4 \\
\bottomrule
\end{tabular}
\end{center}

Se obtiene $\overline d=2.375$ y $s_D=1.685$. Con $\alpha=0.05$, ¿el entrenamiento reduce el tiempo medio de ejecución?

\textbf{Hipótesis:} $H_0:\mu_D\le0$ vs. $H_1:\mu_D>0$ (con $D=$ Antes $-$ Después; si el entrenamiento reduce el tiempo, $D$ debería ser positivo en promedio).

\textbf{Región crítica} (gl$=7$): $T>t_{7,0.05}=1.895$.

\textbf{Cálculo:} $\displaystyle t=\frac{2.375-0}{1.685/\sqrt 8}=\frac{2.375}{0.5958}=3.99$.

\textbf{Valor p:} de tablas de $t_7$, $t_{7,0.005}=3.499$; como $3.99>3.499$, $p<0.005$, muy por debajo de $\alpha=0.05$.

\textbf{Decisión y conclusión:} $3.99>1.895\Rightarrow$ se rechaza $H_0$: el entrenamiento reduce significativamente el tiempo medio de ejecución.
\end{ejemplo}

\subsection{Igualdad de dos varianzas: prueba F de Snedecor}

Sean $X_1,\dots,X_{n_1}\sim N(\mu_1,\sigma_1^2)$ y $Y_1,\dots,Y_{n_2}\sim N(\mu_2,\sigma_2^2)$ muestras independientes. Se desea probar
\[
H_0:\sigma_1^2=\sigma_2^2.
\]
Esta prueba es importante no solo por sí misma (comparar variabilidades es, muchas veces, la pregunta sustantiva del problema), sino porque, como se mencionó, es el criterio formal para decidir si es razonable aplicar el estadístico $t$ agrupado de la subsección anterior.

\begin{teorema}
Bajo $H_0$, el estadístico
\[
F=\frac{S_1^2}{S_2^2}
\]
se distribuye según una $F$ de Snedecor con $n_1-1$ grados de libertad en el numerador y $n_2-1$ en el denominador, $F_{n_1-1,\,n_2-1}$.
\end{teorema}

\begin{proof}
Recordemos que la distribución $F$ con $k_1$ y $k_2$ grados de libertad se define como el cociente de dos variables $\chi^2$ independientes, cada una dividida por sus grados de libertad: si $W_1\sim\chi^2_{k_1}$ y $W_2\sim\chi^2_{k_2}$ son independientes, entonces $\dfrac{W_1/k_1}{W_2/k_2}\sim F_{k_1,k_2}$.

Sabemos que $(n_1-1)S_1^2/\sigma_1^2\sim\chi^2_{n_1-1}$ y $(n_2-1)S_2^2/\sigma_2^2\sim\chi^2_{n_2-1}$, y que ambas son independientes (provienen de muestras independientes). Aplicando la definición con $W_1=(n_1-1)S_1^2/\sigma_1^2$, $k_1=n_1-1$, $W_2=(n_2-1)S_2^2/\sigma_2^2$, $k_2=n_2-1$:
\[
\frac{W_1/k_1}{W_2/k_2} = \frac{\left(\dfrac{(n_1-1)S_1^2/\sigma_1^2}{n_1-1}\right)}{\left(\dfrac{(n_2-1)S_2^2/\sigma_2^2}{n_2-1}\right)}=\frac{S_1^2/\sigma_1^2}{S_2^2/\sigma_2^2}=\frac{S_1^2}{S_2^2}\cdot\frac{\sigma_2^2}{\sigma_1^2}\ \sim\ F_{n_1-1,n_2-1}.
\]
Bajo $H_0:\sigma_1^2=\sigma_2^2$, el factor $\sigma_2^2/\sigma_1^2=1$, de modo que $F=S_1^2/S_2^2\sim F_{n_1-1,n_2-1}$, como se quería demostrar.
\end{proof}

\begin{table}[h]
\centering
\small
\begin{tabular}{@{}lc@{}}
\toprule
$H_1$ & Región crítica \\
\midrule
$\sigma_1^2\neq\sigma_2^2$ & $F<F_{n_1-1,n_2-1,\,1-\alpha/2}$ \ o \ $F>F_{n_1-1,n_2-1,\,\alpha/2}$ \\
$\sigma_1^2>\sigma_2^2$ & $F>F_{n_1-1,n_2-1,\,\alpha}$ \\
\bottomrule
\end{tabular}
\end{table}

\begin{observacion}[Convención práctica y relación con las tablas]
Las tablas usuales de la distribución $F$ solo tabulan cuantiles en la cola superior ($\gamma$ pequeño). Para obtener cuantiles de la cola inferior se utiliza la identidad
\[
F_{k_1,k_2,\,1-\gamma}=\frac{1}{F_{k_2,k_1,\,\gamma}},
\]
que se deduce de que si $F\sim F_{k_1,k_2}$, entonces $1/F\sim F_{k_2,k_1}$ (basta invertir el cociente de las dos $\chi^2$ normalizadas en la definición). Por esta razón, en la práctica —y en particular para el caso bilateral—, es usual colocar en el numerador la muestra con mayor varianza muestral, de modo que el estadístico calculado resulte $F\ge1$, y compararlo únicamente con el valor crítico superior $F_{n_1-1,n_2-1,\,\alpha/2}$ (con los grados de libertad correspondientes a la muestra que efectivamente se colocó en el numerador), evitando así tener que buscar cuantiles de cola inferior.
\end{observacion}

\begin{center}
\begin{tikzpicture}
\begin{axis}[
  width=10cm, height=5cm,
  axis lines=left, ytick=\empty,
  xlabel={$F$}, ylabel={densidad},
  xmin=0, xmax=5, ymin=0,
  legend pos=north east, legend style={font=\footnotesize},
]
\addplot[domain=0.05:5, samples=150, thick, blue] {(x^1)/((1+2*x)^3)};
\addlegendentry{$F_{4,10}$}
\addplot[domain=0.05:5, samples=150, thick, red, dashed] {(x^5)/((1+1.2*x)^8)*3};
\addlegendentry{$F_{12,20}$ (más concentrada en $1$)}
\end{axis}
\end{tikzpicture}
\end{center}

\begin{ejemplo}
Un nuevo instrumento de medición se propone como más preciso (menor variabilidad) que el instrumento estándar. Instrumento estándar: $n_1=16$, $s_1^2=2.8$. Instrumento nuevo: $n_2=13$, $s_2^2=1.1$. Con $\alpha=0.05$, ¿hay evidencia de que el nuevo instrumento es más preciso?

\textbf{1. Hipótesis:} $H_0:\sigma_1^2\le\sigma_2^2$ vs. $H_1:\sigma_1^2>\sigma_2^2$ (equivalente a trabajar con $H_0:\sigma_1^2=\sigma_2^2$).

\textbf{2. Estadístico:} $\displaystyle F=\frac{S_1^2}{S_2^2}\overset{H_0}{\sim}F_{15,12}$.

\textbf{3. Región crítica:} $F>F_{15,12,\,0.05}=2.62$.

\textbf{4. Cálculo:} $\displaystyle F=\frac{2.8}{1.1}=2.55$.

\textbf{5. Decisión:} $2.55<2.62\Rightarrow$ no se rechaza $H_0$.

\textbf{6. Conclusión:} al $5\%$, no hay evidencia estadísticamente suficiente de que el nuevo instrumento sea más preciso, aunque el resultado es sugestivo (muy cercano al valor crítico) y amerita ampliar la muestra antes de descartar la mejora.
\end{ejemplo}

\begin{ejemplo}[Caso bilateral]
Se comparan dos líneas de producción respecto de la variabilidad en el diámetro (mm) de una pieza. Línea A: $n_1=13$, $s_1^2=0.048$. Línea B: $n_2=10$, $s_2^2=0.019$. Con $\alpha=0.10$, ¿hay evidencia de que las variabilidades difieren?

\textbf{1. Hipótesis:} $H_0:\sigma_1^2=\sigma_2^2$ vs. $H_1:\sigma_1^2\neq\sigma_2^2$.

\textbf{2. Estadístico} (numerador = mayor varianza muestral, según la convención práctica): $\displaystyle F=\frac{S_1^2}{S_2^2}\overset{H_0}{\sim}F_{12,9}$.

\textbf{3. Región crítica:} $F>F_{12,9,\,0.05}=3.07$.

\textbf{4. Cálculo:}
\[
F=\frac{0.048}{0.019}=2.53.
\]

\textbf{5. Decisión:} $2.53<3.07\Rightarrow$ no se rechaza $H_0$.

\textbf{6. Conclusión:} al nivel $10\%$, no hay evidencia suficiente de que las variabilidades de ambas líneas difieran, por lo cual, de ser necesario, sería razonable aplicar el estadístico $t$ agrupado (pooled) para comparar las medias de ambas líneas.
\end{ejemplo}

\newpage
\section{Pruebas chi-cuadrado: bondad de ajuste, independencia y homogeneidad}

Las pruebas desarrolladas en las secciones anteriores se aplican a parámetros de variables cuantitativas (medias, varianzas). En esta sección estudiamos una familia distinta de pruebas, basada en la comparación de \textbf{frecuencias observadas} contra \textbf{frecuencias esperadas} bajo una hipótesis nula, aplicables fundamentalmente a datos categóricos (de conteo). Las tres pruebas que veremos —bondad de ajuste, independencia y homogeneidad— comparten el mismo tipo de estadístico, propuesto originalmente por Karl Pearson.

\begin{definicion}[Estadístico $\chi^2$ de Pearson]
\[
\chi^2=\sum_i\frac{(O_i-E_i)^2}{E_i},
\]
donde $O_i$ son las frecuencias \textbf{observadas} y $E_i$ las frecuencias \textbf{esperadas} bajo $H_0$, en cada una de las categorías o celdas consideradas.
\end{definicion}

Intuitivamente, cada término $(O_i-E_i)^2/E_i$ mide, en una escala relativa al tamaño esperado, cuánto se aparta la frecuencia observada en la categoría $i$ de la que predice $H_0$; la división por $E_i$ evita que las categorías con mayor frecuencia esperada dominen artificialmente la suma solo por tener un valor absoluto de discrepancia más grande. Cuando $H_0$ es verdadera, cada término tiende a ser pequeño (en un sentido probabilístico preciso); cuando $H_0$ es falsa, algunas frecuencias observadas se apartan sistemáticamente de las esperadas y el estadístico tiende a tomar valores grandes. Por esta razón, en las tres pruebas de esta sección la región crítica está \textbf{siempre en la cola derecha}: valores grandes de $\chi^2$ (mucha discrepancia) llevan a rechazar $H_0$; valores pequeños de $\chi^2$ (buen ajuste entre lo observado y lo esperado) son compatibles con $H_0$.

\begin{center}
\begin{tikzpicture}
\begin{axis}[
  width=10cm, height=5cm,
  axis lines=left, ytick=\empty,
  xlabel={$\chi^2$}, ylabel={densidad},
  xmin=0, xmax=25, ymin=0,
]
\addplot[domain=0.1:25, samples=150, thick, blue] {x^(1.5)*exp(-x/2)};
\addplot[domain=11.07:25, samples=60, fill=red!40, draw=none] {x^(1.5)*exp(-x/2)} \closedcycle;
\draw[dashed] (axis cs:11.07,0) -- (axis cs:11.07,0.85);
\node at (axis cs:11.07,-0.15) {\footnotesize $\chi^2_{k-1,\alpha}$};
\end{axis}
\end{tikzpicture}
\end{center}

La justificación teórica de por qué este estadístico se distribuye aproximadamente $\chi^2$ bajo $H_0$ proviene de un resultado asintótico: cada frecuencia observada $O_i$ puede pensarse como una suma de indicadores independientes, de modo que, por el Teorema Central del Límite (multivariado), el vector de frecuencias observadas se aproxima a una distribución normal multivariada, y el estadístico de Pearson resulta ser, asintóticamente, una forma cuadrática de esas variables normales que se distribuye $\chi^2$ con tantos grados de libertad como dimensiones libres tenga el problema (el número de categorías menos el número de restricciones impuestas, entre ellas la de que las frecuencias esperadas sumen $n$, y menos los parámetros estimados a partir de los datos).

\subsection{Prueba de bondad de ajuste}

Se dispone de $n$ observaciones clasificadas en $k$ categorías mutuamente excluyentes y exhaustivas, con frecuencias observadas $O_1,\dots,O_k$ (con $\sum_i O_i=n$). Se desea comprobar si los datos son compatibles con una distribución teórica que asigna probabilidades $p_1,\dots,p_k$ a cada categoría ($\sum_i p_i=1$):
\[
H_0:\ P(\text{categoría } i)=p_i,\ i=1,\dots,k \qquad \text{vs.}\qquad H_1:\ \text{para algún } i,\ P(\text{categoría } i)\neq p_i.
\]

Bajo $H_0$, la frecuencia \textbf{esperada} en la categoría $i$ es $E_i=n\,p_i$ (el número medio de observaciones que caerían en esa categoría si $H_0$ fuera cierta), y el estadístico de prueba es
\[
\chi^2=\sum_{i=1}^k\frac{(O_i-E_i)^2}{E_i}\ \overset{H_0}{\approx}\ \chi^2_{k-1-m},
\]
donde $m$ es la cantidad de parámetros de la distribución teórica que debieron \textbf{estimarse} a partir de la propia muestra (si la distribución está completamente especificada de antemano, sin parámetros libres, $m=0$).

\begin{observacion}[Por qué se pierden grados de libertad al estimar parámetros]
Los grados de libertad de referencia son $k-1$ (y no $k$) porque las frecuencias esperadas $E_i=np_i$ están sujetas a la restricción $\sum_i E_i=n=\sum_i O_i$: solo $k-1$ de las $k$ discrepancias $(O_i-E_i)$ son libres, la última queda determinada por las demás. Cuando, además, algún parámetro de la distribución teórica (por ejemplo, la media $\lambda$ de una Poisson, o la media y varianza de una normal) debe \textbf{estimarse} a partir de los mismos datos $O_1,\dots,O_k$ que se están usando para calcular el estadístico, se introduce una restricción adicional por cada parámetro estimado: intuitivamente, los valores estimados "se ajustan" a los datos observados, acercando artificialmente las frecuencias esperadas a las observadas y reduciendo la discrepancia esperable bajo $H_0$. Este efecto se compensa restando un grado de libertad adicional por cada parámetro estimado, lo que da la fórmula general $k-1-m$.
\end{observacion}

\begin{observacion}[Condición de aplicación]
La aproximación de la distribución exacta (discreta) del estadístico $\chi^2$ de Pearson por la distribución continua $\chi^2_{k-1-m}$ es adecuada únicamente cuando las frecuencias esperadas no son demasiado pequeñas. El criterio usual —muchas veces llamado \textbf{condición de Cochran}— exige
\[
E_i\ge5 \quad \text{para todas las categorías (o, en versiones más flexibles, para al menos el $80\%$ de ellas, sin ninguna celda con $E_i<1$).}
\]
Si alguna $E_i$ no cumple esta condición, se recomienda \textbf{agrupar} categorías contiguas (o categorías con sentido similar) hasta satisfacerla, teniendo en cuenta que esto reduce en consecuencia el número efectivo de categorías $k$ y, por lo tanto, los grados de libertad de la prueba.
\end{observacion}

\begin{ejemplo}[Distribución completamente especificada, $m=0$]
Se lanza un dado $120$ veces para verificar si está equilibrado. Frecuencias observadas:
\begin{center}
\begin{tabular}{@{}lcccccc@{}}
\toprule
Cara & 1 & 2 & 3 & 4 & 5 & 6 \\
\midrule
$O_i$ & 25 & 17 & 15 & 23 & 24 & 16 \\
\bottomrule
\end{tabular}
\end{center}
Con $\alpha=0.05$, pruebe si el dado es equilibrado.

\textbf{1. Hipótesis:} $H_0:p_i=1/6\ \forall i$ vs. $H_1:$ algún $p_i\neq1/6$.

\textbf{2. Frecuencias esperadas:} $E_i=120\times\tfrac16=20$ para todo $i$ (todas $\ge5$: se cumple la condición de Cochran).

\textbf{3. Grados de libertad:} $k=6$, $m=0\Rightarrow \text{gl}=5$.

\textbf{4. Región crítica:} $\chi^2>\chi^2_{5,0.05}=11.07$.

\textbf{5. Cálculo:}
\[
\chi^2=\frac{5^2}{20}+\frac{3^2}{20}+\frac{5^2}{20}+\frac{3^2}{20}+\frac{4^2}{20}+\frac{4^2}{20}=\frac{25+9+25+9+16+16}{20}=\frac{100}{20}=5.0.
\]

\textbf{6. Decisión y conclusión:} $5.0<11.07\Rightarrow$ no se rechaza $H_0$; no hay evidencia de que el dado esté desequilibrado. (De tablas de $\chi^2_5$, $\chi^2_{5,0.50}\approx4.35$, de modo que el valor p es algo mayor que $0.50$: la discrepancia observada es incluso menor que la típica bajo $H_0$.)
\end{ejemplo}

\begin{ejemplo}[Distribución con parámetro estimado, $m=1$]
Se registra el número de defectos por unidad en $200$ piezas, y se desea comprobar si sigue una distribución de Poisson (la media $\lambda$ debe estimarse de los datos: $\hat\lambda=1.5$).
\begin{center}
\small
\begin{tabular}{@{}lccccc@{}}
\toprule
N° defectos & 0 & 1 & 2 & 3 & $\ge4$ \\
\midrule
$O_i$ & 48 & 65 & 48 & 26 & 13 \\
$p_i$ (Poisson, $\hat\lambda=1.5$) & 0.223 & 0.335 & 0.251 & 0.126 & 0.065 \\
$E_i=200\,p_i$ & 44.6 & 67.0 & 50.2 & 25.2 & 13.0 \\
\bottomrule
\end{tabular}
\end{center}
Con $\alpha=0.05$: gl$=k-1-m=5-1-1=3$ (se estimó $\lambda$, por lo que $m=1$).

\textbf{Región crítica:} $\chi^2>\chi^2_{3,0.05}=7.815$.

\textbf{Cálculo:}
\[
\chi^2=\frac{(48-44.6)^2}{44.6}+\frac{(65-67.0)^2}{67.0}+\frac{(48-50.2)^2}{50.2}+\frac{(26-25.2)^2}{25.2}+\frac{(13-13.0)^2}{13.0}
\]
\[
\approx 0.259+0.060+0.096+0.025+0.000=0.44.
\]

\textbf{Decisión:} $0.44<7.815\Rightarrow$ no se rechaza $H_0$: el ajuste a la distribución de Poisson es adecuado. El valor p es muy alto (bastante mayor que $0.90$), consistente con un ajuste casi perfecto entre lo observado y lo esperado.
\end{ejemplo}

\subsection{Prueba de independencia en tablas de contingencia}

Se clasifica una muestra de $n$ individuos simultáneamente según \textbf{dos variables categóricas}, $A$ (con $r$ categorías) y $B$ (con $c$ categorías), formando una tabla de contingencia de $r\times c$ celdas con frecuencias observadas $O_{ij}$, $i=1,\dots,r$, $j=1,\dots,c$. Denotamos $O_{i\bullet}=\sum_j O_{ij}$ el total de la fila $i$ y $O_{\bullet j}=\sum_i O_{ij}$ el total de la columna $j$, con $\sum_i O_{i\bullet}=\sum_j O_{\bullet j}=n$.

\[
H_0:\ A \text{ y } B \text{ son independientes} \qquad \text{vs.}\qquad H_1:\ A \text{ y } B \text{ no son independientes}.
\]

\begin{teorema}[Frecuencias esperadas bajo independencia]
Bajo el supuesto de independencia entre $A$ y $B$, la frecuencia esperada en la celda $(i,j)$ es
\[
E_{ij}=\frac{O_{i\bullet}\cdot O_{\bullet j}}{n}.
\]
\end{teorema}

\begin{proof}
Sean $p_{i\bullet}=P(A=i)$ y $p_{\bullet j}=P(B=j)$ las probabilidades marginales (desconocidas) de cada categoría. Si $A$ y $B$ son independientes, la probabilidad conjunta de la celda $(i,j)$ es $p_{ij}=p_{i\bullet}\cdot p_{\bullet j}$ (definición de independencia de eventos), y la frecuencia esperada en esa celda, sobre $n$ observaciones, es $E_{ij}=n\,p_{ij}=n\,p_{i\bullet}\,p_{\bullet j}$.

Como las probabilidades marginales $p_{i\bullet}$ y $p_{\bullet j}$ son desconocidas, se estiman naturalmente por sus proporciones muestrales, $\hat p_{i\bullet}=O_{i\bullet}/n$ y $\hat p_{\bullet j}=O_{\bullet j}/n$ (estimadores de máxima verosimilitud de las marginales). Sustituyendo estas estimaciones en la expresión de $E_{ij}$:
\[
\widehat E_{ij}=n\cdot\frac{O_{i\bullet}}{n}\cdot\frac{O_{\bullet j}}{n}=\frac{O_{i\bullet}\cdot O_{\bullet j}}{n},
\]
que es la fórmula enunciada.
\end{proof}

El estadístico de prueba es
\[
\chi^2=\sum_{i=1}^r\sum_{j=1}^c \frac{(O_{ij}-E_{ij})^2}{E_{ij}}\ \overset{H_0}{\approx}\ \chi^2_{(r-1)(c-1)}.
\]

\begin{observacion}[Grados de libertad]
Los $(r-1)(c-1)$ grados de libertad se explican de la siguiente manera: la tabla tiene $rc$ celdas, pero las frecuencias esperadas $E_{ij}$ se calculan a partir de las $r+c-1$ marginales estimadas de la muestra (los $r$ totales de fila y los $c$ totales de columna, que suman una restricción menos porque ambos conjuntos de totales deben sumar $n$). Restando estas restricciones —y compensando adecuadamente la doble estimación—, el número de celdas cuya discrepancia $(O_{ij}-E_{ij})$ es efectivamente libre resulta $rc-(r-1)-(c-1)-1=(r-1)(c-1)$.
\end{observacion}

\begin{ejemplo}
Se estudia si el ausentismo laboral es independiente del turno de trabajo, sobre $300$ empleados:
\begin{center}
\begin{tabular}{@{}lccc|c@{}}
\toprule
 & Mañana & Tarde & Noche & Total \\
\midrule
Con ausentismo & 18 & 22 & 35 & 75 \\
Sin ausentismo & 82 & 78 & 65 & 225 \\
\midrule
Total & 100 & 100 & 100 & 300 \\
\bottomrule
\end{tabular}
\end{center}
Con $\alpha=0.05$, ¿el ausentismo depende del turno?

\textbf{1. Hipótesis:} $H_0:$ ausentismo y turno son independientes vs. $H_1:$ no lo son.

\textbf{2. Frecuencias esperadas} ($E_{ij}=O_{i\bullet}O_{\bullet j}/300$):
\[
E_{11}=\frac{75\times100}{300}=25,\quad E_{12}=25,\quad E_{13}=25,\quad E_{21}=75,\quad E_{22}=75,\quad E_{23}=75.
\]
Todas las frecuencias esperadas superan $5$: se cumple la condición de Cochran.

\textbf{3. Grados de libertad:} $(r-1)(c-1)=(2-1)(3-1)=2$.

\textbf{4. Región crítica:} $\chi^2>\chi^2_{2,0.05}=5.991$.

\textbf{5. Cálculo:}
\[
\chi^2=\frac{(18-25)^2}{25}+\frac{(22-25)^2}{25}+\frac{(35-25)^2}{25}+\frac{(82-75)^2}{75}+\frac{(78-75)^2}{75}+\frac{(65-75)^2}{75}
\]
\[
=\frac{49}{25}+\frac{9}{25}+\frac{100}{25}+\frac{49}{75}+\frac{9}{75}+\frac{100}{75}=1.96+0.36+4.00+0.653+0.12+1.333=8.43.
\]

\textbf{6. Decisión:} $8.43>5.991\Rightarrow$ se rechaza $H_0$.

\textbf{7. Valor p:} de tablas de $\chi^2_2$, $\chi^2_{2,0.025}=7.378$ y $\chi^2_{2,0.01}=9.210$; como $7.378<8.43<9.210$, el valor p está entre $0.01$ y $0.025$.

\textbf{8. Conclusión:} al $5\%$ de significación, hay evidencia de que el ausentismo \textbf{no} es independiente del turno de trabajo; el turno noche muestra un ausentismo notablemente mayor que el esperado bajo independencia.
\end{ejemplo}

\subsection{Prueba de homogeneidad}

Se toman muestras \textbf{independientes} de $r$ poblaciones distintas, con tamaños fijados de antemano por el diseño del muestreo, y se clasifica cada observación según una variable categórica común con $c$ categorías. Se desea probar si las $r$ poblaciones tienen la \textbf{misma distribución} respecto de esa variable:
\[
H_0:\ \text{las } r \text{ poblaciones tienen la misma distribución en la variable categórica},
\]
\[
H_1:\ \text{al menos una población difiere de las demás}.
\]

\begin{observacion}[Diferencia conceptual con la prueba de independencia]
La distinción entre independencia y homogeneidad es puramente de \textbf{diseño de muestreo}, no de cálculo. En la prueba de \textbf{independencia}, se extrae \textbf{una única muestra aleatoria} de $n$ individuos de una población, y se los clasifica \emph{a posteriori} según dos variables categóricas; en ese caso, tanto los totales de fila como los de columna son cantidades \textbf{aleatorias}, resultado del muestreo. En la prueba de \textbf{homogeneidad}, en cambio, se extraen $r$ muestras \textbf{independientes}, una de cada población, con tamaños (los totales de fila, o de columna, según cómo se organice la tabla) \textbf{fijados de antemano} por el investigador, y la variable de clasificación es una sola (la que define las columnas).

A pesar de esta diferencia conceptual en el mecanismo de generación de los datos, el desarrollo matemático —frecuencias esperadas $E_{ij}=O_{i\bullet}O_{\bullet j}/n$, estadístico $\chi^2=\sum (O_{ij}-E_{ij})^2/E_{ij}$, y grados de libertad $(r-1)(c-1)$— resulta \textbf{formalmente idéntico} al de la prueba de independencia. La razón profunda es que, condicionalmente a los totales marginales observados, la distribución del estadístico bajo ambas hipótesis nulas (independencia, homogeneidad) coincide.
\end{observacion}

\begin{ejemplo}
Se extraen muestras de $150$ piezas de cada uno de tres proveedores y se clasifican según su calidad:
\begin{center}
\begin{tabular}{@{}lccc|c@{}}
\toprule
 & Proveedor A & Proveedor B & Proveedor C & Total \\
\midrule
Primera calidad & 120 & 108 & 96 & 324 \\
Segunda calidad & 30 & 42 & 54 & 126 \\
\midrule
Total & 150 & 150 & 150 & 450 \\
\bottomrule
\end{tabular}
\end{center}
Con $\alpha=0.01$, ¿los tres proveedores son homogéneos respecto de la calidad entregada?

\textbf{1. Hipótesis:} $H_0:$ los tres proveedores tienen igual distribución de calidad vs. $H_1:$ al menos uno difiere.

\textbf{2. Frecuencias esperadas:} $E_{1j}=\dfrac{324\times150}{450}=108$ para cada proveedor (primera calidad); $E_{2j}=\dfrac{126\times150}{450}=42$ (segunda calidad). Todas superan $5$: se cumple la condición de Cochran.

\textbf{3. Grados de libertad:} $(r-1)(c-1)=(2-1)(3-1)=2$.

\textbf{4. Región crítica:} $\chi^2>\chi^2_{2,0.01}=9.210$.

\textbf{5. Cálculo:}
\[
\chi^2=\frac{(120-108)^2}{108}+\frac{(108-108)^2}{108}+\frac{(96-108)^2}{108}+\frac{(30-42)^2}{42}+\frac{(42-42)^2}{42}+\frac{(54-42)^2}{42}
\]
\[
=\frac{144}{108}+0+\frac{144}{108}+\frac{144}{42}+0+\frac{144}{42}=1.333+1.333+3.429+3.429=9.52.
\]

\textbf{6. Decisión:} $9.52>9.210\Rightarrow$ se rechaza $H_0$.

\textbf{7. Valor p:} de tablas de $\chi^2_2$, $\chi^2_{2,0.01}=9.210$ y $\chi^2_{2,0.005}=10.60$; como $9.210<9.52<10.60$, el valor p está entre $0.005$ y $0.01$.

\textbf{8. Conclusión:} al $1\%$ de significación, hay evidencia de que los tres proveedores no son homogéneos en la calidad de las piezas entregadas; el proveedor C presenta una proporción de segunda calidad notablemente mayor que la esperada bajo homogeneidad.
\end{ejemplo}

\subsection{Comparación de las tres pruebas y condiciones de aplicación}

\begin{table}[h]
\centering
\small
\begin{tabular}{@{}p{2.7cm}p{4.3cm}p{5.3cm}@{}}
\toprule
Prueba & Diseño de muestreo & Pregunta que responde \\
\midrule
Bondad de ajuste & una muestra, una variable categórica & ¿la variable sigue una distribución teórica dada? \\[0.6em]
Independencia & una muestra, dos variables categóricas (ambos totales marginales aleatorios) & ¿las dos variables están relacionadas? \\[0.6em]
Homogeneidad & varias muestras independientes (una por población, tamaños fijados de antemano), una variable categórica & ¿las poblaciones tienen la misma distribución? \\
\bottomrule
\end{tabular}
\caption{Comparación de las tres pruebas basadas en el estadístico de Pearson.}
\end{table}

Pese a responder preguntas conceptualmente distintas, las pruebas de independencia y homogeneidad comparten exactamente el mismo cálculo. Además, todas las pruebas $\chi^2$ desarrolladas en esta sección están sujetas a las mismas condiciones generales de aplicación:

\begin{itemize}[nosep]
\item Las observaciones deben ser \textbf{independientes} entre sí.
\item Cada individuo u observación debe clasificarse en \textbf{una y solo una} categoría o celda (categorías mutuamente excluyentes y exhaustivas).
\item Las frecuencias esperadas deben ser suficientemente grandes para que la aproximación por la distribución continua $\chi^2$ sea adecuada: la \textbf{condición de Cochran} exige $E_i\ge5$ (o $E_{ij}\ge5$ en tablas de contingencia) en todas las celdas, admitiéndose en versiones más flexibles hasta un $20\%$ de celdas con $E_i$ entre $1$ y $5$, siempre que ninguna celda tenga $E_i<1$.
\item Si la condición no se cumple, se deben \textbf{agrupar} categorías o filas/columnas adyacentes (con sentido sustantivo) hasta satisfacerla, lo cual reduce en consecuencia los grados de libertad de la prueba.
\item En las tres pruebas, el estadístico de Pearson conduce siempre a una región crítica de \textbf{cola derecha}: valores grandes de $\chi^2$ indican una discrepancia importante entre lo observado y lo esperado bajo $H_0$, y por lo tanto evidencia en contra de $H_0$.
\end{itemize}

\newpage
\section{Observaciones adicionales}

A modo de cierre, conviene situar el contenido de esta unidad dentro del panorama más amplio de la inferencia estadística, y anticipar su conexión con la Unidad 8.

\begin{itemize}[nosep]
\item \textbf{Unificación conceptual.} A lo largo de esta unidad vimos una variedad de estadísticos de prueba —$Z$, $T$, $\chi^2$, $F$— aplicados a distintos parámetros y situaciones. Es importante no perder de vista que, detrás de esa aparente diversidad, subyace siempre la misma estructura lógica: un estadístico cuya distribución bajo $H_0$ es conocida (exacta o aproximadamente), una región crítica determinada por $\alpha$ y por el sentido de $H_1$, y una decisión basada en comparar lo observado con lo esperado bajo $H_0$, ya sea mediante el estadístico o mediante el valor p.
\item \textbf{Relación con la estimación.} La dualidad entre pruebas de hipótesis e intervalos de confianza, demostrada en la Sección 1, no es una curiosidad aislada: refleja que estimación y prueba de hipótesis son, en el fondo, dos maneras de resumir la misma información muestral (a través del mismo estadístico pivotal) para responder preguntas distintas —"¿cuál es el rango plausible de $\theta$?" versus "¿es $\theta$ igual a un valor específico?"—.
\item \textbf{Hacia la Unidad 8: el modelo de regresión lineal.} En la Unidad 8 se introduce el modelo de regresión lineal simple, $Y_i=\beta_0+\beta_1 x_i+\varepsilon_i$, con $\varepsilon_i\sim N(0,\sigma^2)$ independientes. Bajo este modelo, los estimadores de mínimos cuadrados $\hat\beta_0$ y $\hat\beta_1$ resultan ser combinaciones lineales de las observaciones $Y_i$, y por lo tanto normales; esto permite construir, exactamente con la misma lógica que la prueba $t$ para una media desarrollada en la Sección 2, estadísticos de la forma
\[
T=\frac{\hat\beta_1-\beta_{1,0}}{S_{\hat\beta_1}}\ \sim\ t_{n-2}
\]
para probar hipótesis sobre la pendiente poblacional (típicamente $H_0:\beta_1=0$, es decir, ausencia de relación lineal entre $x$ e $y$). La construcción de este estadístico —numerador que mide un apartamiento respecto de un valor de referencia, denominador que estima el error estándar correspondiente, cociente que resulta $t$ de Student— es exactamente la misma estructura que la prueba $t$ para una media vista en este apunte, aplicada ahora a un parámetro de un modelo más complejo.
\item \textbf{Hacia la Unidad 8: el análisis de varianza (ANOVA).} La prueba $F$ de Snedecor, introducida en la Sección 3 para comparar dos varianzas, reaparece en la Unidad 8 en un contexto distinto pero estrechamente emparentado: el análisis de varianza de un factor (ANOVA), utilizado para comparar las medias de \textbf{más de dos} poblaciones simultáneamente (una generalización natural de la comparación de dos medias vista en la Sección 3 de este apunte). La idea del ANOVA es descomponer la variabilidad total de los datos en una componente \textbf{entre grupos} y una componente \textbf{dentro de los grupos}, y comparar ambas mediante un cociente de varianzas —formalmente análogo al estadístico $F=S_1^2/S_2^2$ ya estudiado— que se distribuye $F$ bajo la hipótesis nula de igualdad de todas las medias poblacionales. Es decir, el mismo estadístico $F$ que aquí usamos para contrastar variabilidades se reutiliza, con una interpretación distinta de "numerador" y "denominador", para contrastar medias de múltiples grupos.
\item \textbf{Una perspectiva unificadora.} Tanto la regresión como el ANOVA pueden presentarse como \emph{modelos lineales}, y las pruebas de hipótesis asociadas a sus parámetros (pendientes, medias de grupo) no son más que aplicaciones —más elaboradas, pero conceptualmente idénticas— del marco de prueba de hipótesis desarrollado íntegramente en esta unidad: formular $H_0$ y $H_1$, elegir un estadístico con distribución conocida bajo $H_0$, determinar la región crítica o el valor p, y decidir. Dominar en profundidad los conceptos de esta Unidad 7 —en particular, la lógica de los errores tipo I y II, la construcción de estadísticos a partir de pivotes normales y $\chi^2$, y la interpretación correcta del valor p— es, por lo tanto, el requisito indispensable para abordar con solidez los modelos más ricos que se estudian en la Unidad 8.
\end{itemize}

\end{document}
